\documentclass[12pt]{article}

\usepackage[margin=1in]{geometry}
\usepackage{fancyhdr}
\usepackage{sectsty}
\usepackage{titlesec}
\usepackage{rotating}
\usepackage{lipsum}
\usepackage{setspace}
\usepackage{enumitem}
\usepackage{authblk}
\usepackage{etoolbox}
\usepackage{changepage}
\usepackage{xpatch}
\usepackage{threeparttable}
\usepackage{array}
\usepackage{booktabs}
% POQ table style: every rule (top, header, spanner, bottom) the same thin weight
\setlength{\heavyrulewidth}{0.5pt}
\setlength{\lightrulewidth}{0.5pt}
\setlength{\cmidrulewidth}{0.5pt}
\usepackage{tabularx}
\usepackage{marvosym}
\usepackage[title]{appendix}
\usepackage{caption}
\usepackage{subcaption}
\usepackage{wrapfig}
\usepackage{graphicx}
\graphicspath{ {figures/} }
\usepackage{epigraph}
\usepackage{titling}
\usepackage{blindtext}
\usepackage{float}
\usepackage{amsmath,amssymb}
\usepackage{pdflscape}
\usepackage{multirow}
\usepackage{adjustbox}
\usepackage{xcolor}
\usepackage[table]{xcolor}
\usepackage{lmodern}
\usepackage{microtype}
\usepackage{chngcntr}

% --------------------------------------------------
% Font & spacing
% --------------------------------------------------
\renewcommand{\familydefault}{\rmdefault}
\doublespacing

% --------------------------------------------------
% Header / footer
% --------------------------------------------------
\pagestyle{fancy}
\fancyhf{}
\fancyfoot[R]{\thepage}
\renewcommand{\headrulewidth}{0pt} % no rule at the top of each page

% Title page style (no headers or footers)
\fancypagestyle{titlepage}{
  \fancyhf{}
}

% Horizontal rule after title
\newcommand{\HRule}{\noindent\rule{\linewidth}{0.5pt}}

% Title, set flush left and bold as in POQ
\newcommand{\papertitle}{The Accuracy of Matched Sample Surveys: Evidence from the Cooperative Election Study}
\newcommand{\printtitle}{%
  \begin{singlespace}\noindent{\raggedright\bfseries\Large\papertitle\par}\end{singlespace}}

% Abstract envi
\newenvironment{myabstract}
  {%
    \begin{adjustwidth}{4em}{0em} % left indent = 2em, right indent = 0em
    \noindent\textbf{Abstract}\hspace{1em}%
  }
  {%
    \end{adjustwidth}
    \par
  }

% Table styling
\definecolor{primarycell}{RGB}{255,235,205}
  
% Author footnote numbers
\renewcommand\Authsep{, }
\renewcommand\Authands{, }
\renewcommand\Affilfont{\small}

% --------------------------------------------------
% Section headings (POQ style)
% --------------------------------------------------
% Level 1: centered, bold. Level 2: flush left, bold. Level 3: flush left,
% italic. Serif throughout, unnumbered.
\titleformat{\section}
  {\normalfont\bfseries\large\centering}
  {}
  {0pt}
  {}

\titleformat{\subsection}
  {\normalfont\bfseries\normalsize\raggedright}
  {}
  {0pt}
  {}

\titleformat{\subsubsection}
  {\normalfont\itshape\normalsize\raggedright}
  {}
  {0pt}
  {}

% Level 4 (used only for the item headings in Supplementary Material
% section F): italic, run in to the paragraph, followed by a period.
\titleformat{\paragraph}[runin]
  {\normalfont\itshape\normalsize}
  {}
  {0pt}
  {}
  [.]
\titlespacing*{\paragraph}{0pt}{1.5ex plus .5ex minus .2ex}{0.5em}

% --------------------------------------------------
% Captions (POQ style)
% --------------------------------------------------
% Bold label followed by a period ("Table 1."), sentence-case text, flush
% left. Table captions sit above the table, figure captions below it.
\captionsetup{labelfont=bf, labelsep=period, justification=raggedright,
  singlelinecheck=false, font=small}
\captionsetup[table]{position=top}

% Bibliography
\usepackage[colorlinks,
            citecolor=blue,   % citations
            linkcolor=blue,   % internal links (e.g. toc)
            urlcolor=blue]{hyperref}
\usepackage{xurl}

% Chicago author-date, the style POQ uses: year after the names, volume:pages,
% no issue numbers, publisher without city, no DOIs.
\usepackage[
  authordate,
  backend=biber,
  compresspages=true,  % page ranges shortened as POQ prints them (307--29)
  natbib=true,          % allow \citep / \citet if you like
  sortcites=true,       % order multiple cites in one parenthesis as in the reference list
  maxcitenames=3,       % in-text: show up to 3 names
  mincitenames=1,       % if more than 3 authors, truncate to 1 + "et al."
  maxbibnames=99,       % in bibliography: list every author
  uniquename=false,     % never add first initials to disambiguate same-surname authors
  uniquelist=false
]{biblatex-chicago}
% let biblatex break long URLs/DOIs at more points
\setcounter{biburlnumpenalty}{7000}
\setcounter{biburlucpenalty}{8000}
\setcounter{biburllcpenalty}{9000}

% make sure multiple keys in one \parencite are separated with "; "
\DeclareDelimFormat{multicitedelim}{\addsemicolon\space}

\AtEveryBibitem{%
  % Print only the year, not months.
  \clearfield{month}%
  \clearfield{day}%
  % Entries with a DOI print no URL either.
  \iffieldundef{doi}{}{\clearfield{url}\clearfield{urlyear}\clearfield{urlmonth}\clearfield{urlday}}%
  % Access dates only for undated online sources (Chicago practice).
  \iffieldundef{year}{}{\clearfield{urlyear}\clearfield{urlmonth}\clearfield{urlday}}%
  % POQ prints no DOIs, ISSNs, or ISBNs.
  \clearfield{doi}\clearfield{issn}\clearfield{isbn}%
  % POQ: journal articles as volume:pages, with no issue number.
  \ifentrytype{article}{\clearfield{number}\clearfield{issue}}{}%
  % POQ: books name the publisher without its city.
  \ifboolexpr{test {\ifentrytype{book}} or test {\ifentrytype{incollection}}
    or test {\ifentrytype{inbook}} or test {\ifentrytype{collection}}}
    {\clearlist{location}}{}%
}

% Some entries store the DOI as a full URL, which then prints the prefix twice.
\DeclareSourcemap{
  \maps[datatype=bibtex]{
    \map{
      \step[fieldsource=doi, match=\regexp{\A\s*https?://(dx\.)?doi\.org/}, replace={}]
    }
  }
}
\addbibresource{CESBib.bib}

% Numbers quoted in the text, captions, and alt text. Generated by the Prose
% Numbers cell of tables_and_figures/figuresHQ.ipynb; do not type them by hand.
\input{output/prose_numbers}


\begin{document}

% ==================================================
% TITLE PAGE
% ==================================================
\begin{titlepage}
\thispagestyle{titlepage}

\printtitle
\vspace{0.4em}
\HRule
\vspace{1.2em}

\begin{singlespace}
\begin{flushleft}
  Jefferson D. Pruett\textsuperscript{1},
  Jon A. Krosnick\textsuperscript{2,*},
  Gabriel Miao Li\textsuperscript{3}\\[0.8em]

  {\footnotesize
    \textsuperscript{1}Department of Economics, Stanford University, Stanford, CA, US\\
    \textsuperscript{2}Professor, Departments of Communication, Political Science, Psychology, and Environmental Social Sciences, Stanford University, Stanford, CA, US\\
    \textsuperscript{3}Assistant Professor, Department of Communication Studies, Chapman University, Orange, CA, US
  }
\end{flushleft}
\end{singlespace}

{\renewcommand{\thefootnote}{*}%
\footnotetext{Corresponding author: Jon A. Krosnick, Stanford University, 432 McClatchy Hall, 450 Jane Stanford Way, Stanford, CA 94305; \texttt{krosnick@stanford.edu}.}}

\vfill

\end{titlepage}

% ==================================================
% MAIN MANUSCRIPT
% ==================================================

\setcounter{page}{1}
\pagestyle{fancy}

\printtitle
\vspace{0.4em}
\HRule
\vspace{1em}

\begin{myabstract}
Nonprobability samples have grown in popularity among survey researchers, with proponents claiming that they can deliver accurate estimates at lower cost than probability samples. This growth has occurred despite accumulating evidence that nonprobability samples are generally less accurate. A notable exception appears in a small number of evaluations of YouGov, particularly a Pew study, which found YouGov to yield more accurate measurements than other nonprobability sample providers. Because YouGov uniquely relies on sample matching, perhaps its accuracy is attributable to that approach. This paper presents the largest evaluation of YouGov data to date, assessing accuracy across nine waves of the Cooperative Election Study, from 2006 to 2022. Across \NObsDemographic{} demographic estimates, \NObsVotingAdmin{} turnout and voting-method estimates, and \NObsCandidate{} candidate-choice estimates, the CES manifested substantial error, often quite large, even when using YouGov’s proprietary weights. Accuracy did not improve over the 17-year period. Without those weights, accuracy was considerably lower, meaning that sample matching alone yielded less accurate measurements. Conventional demographic post-stratification slightly increased average error for the three vote-choice measures observed in every wave and excluded from all weighting, consistent with prior studies.

\end{myabstract}

\clearpage
\printtitle
\vspace{1em}

\section*{Introduction}

Following eye-opening embarrassments for pollsters a century ago stemming from nonprobability sampling, survey researchers have relied on random sampling for decades to produce highly accurate measurements. Yet with the emergence of the internet in 2000, nonprobability sampling reemerged as a dominant methodology. Some firms implement probability sampling for online questionnaire administration. But other firms have collected data more inexpensively by abandoning probability sampling. These providers have often claimed to be implementing sophisticated methods to yield remarkably accurate measurements.

However, many studies have shown that nonprobability
samples yielded less accurate measurements than probability samples (see \cite{cornesse_review_2020} for a review of 25 studies) and in some cases, strikingly less accurate
\parencite[e.g.,][]{macinnis_accuracy_2018, yeager_comparing_2011}. Nonetheless, some scholars
and survey firms have argued that advanced nonprobability methods can yield
accurate estimates \parencite{ansolabehere_cooperative_2013, crampton_about_2007}, raising the
possibility that many prior evaluations did not reflect the performance of more
sophisticated vendors \parencite[see][]{baker_report_2013}.

The present study is a large-scale evaluation of the
accuracy of YouGov’s nonprobability sample surveys, which have used the method
of sample matching \parencite{rivers_sampling_2007}. Sample matching was used initially for causal
inference; matched control groups were constructed
to estimate treatment effects \parencite{stuart_matching_2010}. Applied to internet surveys, this
approach selects a sample from among nonrandom panel members to
approximate characteristics of a probability sample in terms of a small set of
variables, usually demographics \parencite{rivers_sample_2006}. \textcite{noauthor_yougov_nodate} has claimed to be
one of the most quoted pollsters and market research firms in the world. The company
operates in more than 55 countries and is frequently cited by global media.

We focus on nine Cooperative Election Study (CES) surveys conducted over a 17-year period by YouGov, all of which make use of sample matching. The CES, initially called the Cooperative Congressional Election Study (CCES), is the product of a collaboration of universities that contributed funding, supplemented by support from the National Science Foundation. YouGov collected data from large samples before and after each national election in the US since 2006. CES data have been widely used across political science, cited over 4,700 times in scholarly publications.\footnote{Based on a Google Scholar full-text search conducted in August 2026 for ``Cooperative Congressional Election Study'' OR ``Cooperative Election Study'' that returned approximately 4,770 works.}

We begin below by reviewing theoretical views of probability and non-probability sampling, then briefly describe the history of these methods’ use in American survey research during the last century. We then report assessments of the accuracy of CES measurements, broken down to distinguish the error remaining after sample matching from the change produced by standard post-stratification weighting.
% We address these issues using nine Cooperative Election Study (CES) surveys collected over 17 years. We first report accuracy using YouGov’s proprietary weights, focusing on variables not used in weight construction, and then evaluate accuracy without post-survey weighting and with newly constructed calibration weights, isolating the accuracy gains attributable to sample matching versus weighting. The scale and matched-sample design of the CES allow us to assess whether sample matching—and YouGov specifically—constitutes an outlier among nonprobability survey methods and vendors. Throughout, we distinguish primary variables, which were used to generate the CES weights and so are constrained to population levels, from secondary variables, which were not. This distinction is central to interpreting CES accuracy and may help explain the differnece between prior valuations and ours.

\section*{Nonprobability Sampling and the Rise of YouGov}

Probability sampling selects units at random with known inclusion probabilities, ideally so that every member of a population has a known non-zero probability of becoming part of the intended sample. In theory, if members of such a sample are invited to complete a questionnaire, the result should closely resemble what would be obtained if all members of the population had done so.  

% These inferential properties made probability sampling the foundation of major U.S. survey programs, including the Census Bureau’s Current Population Survey (CPS) and American Community Survey (ACS), the Bureau of Labor Statistics’ Consumer Expenditure Surveys, and the CDC’s National Health Interview Survey. After a series of high-profile polling failures in the 1930s and 1940s attributed to nonprobability designs—most notably the Gallup, Crossley, and Roper error in 1948 \parencite{mosteller_pre_election_1949}—probability sampling became the standard for high-quality social scientific research \parencite{cornesse_review_2020}.

When large-scale survey research began in the US, however, probability sampling was not the norm. Well-known researchers, including Gallup, Crossley, and Roper, disseminated results of surveys of non-probability samples, to great acclaim \parencite{crossley_straw_1937}. Only after this methodology yielded embarrassing inaccuracy, especially in 1948 (see \cite{mosteller_pre_election_1949}), did probability sampling become the standard for high-quality social scientific research \parencite{cornesse_review_2020}. This methodology has been the foundation of major US survey programs for decades, including the Census Bureau’s Current Population Survey (CPS) and American Community Survey (ACS), the Bureau of Labor Statistics’ Consumer Expenditure Surveys, the CDC’s National Health Interview Survey, and others. 

With the arrival of the internet, enterprising survey researchers drew probability samples and invited people to complete online questionnaires, increasing efficiency (compared to face-to-face and telephone interviewing) and affording the privacy of self-completion. But firms soon recognized that opt-in nonprobability panels offered even greater speed and cost savings. Although some observers dismissed non-probability samples as “quick and dirty” tools \parencite{bethlehem_selection_2010}, their low per-respondent costs made such web surveys attractive to researchers \parencite{wright_researching_2005}.

As nonprobability samples gained popularity, a substantial benchmarking literature emerged comparing them with high-quality probability samples across domains including voting behavior, consumption, sexual attitudes, and demographics. These studies generally found that probability samples yielded more accurate estimates, whereas nonprobability samples were less accurate and more variable in their accuracy \parencite[e.g.,][]{chang_national_2009, dutwin_apples_2017, erens_nonprobability_2014, macinnis_accuracy_2018, sturgis_assessment_2018, szolnoki_online_2013, yeager_comparing_2011}. Synthesizing evidence from 25 studies, \textcite{cornesse_review_2020} concluded that ``even in the age of declining response rates, accuracy in probability sample surveys is generally higher than in nonprobability sample surveys'' (22).

Researchers began to ask whether post-survey adjustment—especially weighting—can compensate for selection bias in opt-in panels and yield accurate population estimates. A 2013 AAPOR review concluded that such adjustments sometimes reduced bias but produced inconsistent results across outcomes and did not reliably eliminate coverage, nonresponse, and selection biases \parencite{baker_report_2013}.

Methods have evolved, though, and some researchers have suggested that sophisticated techniques can reduce bias. For example, \textcite{wang_forecasting_2015} used multilevel regression and post-stratification to obtain accurate 2012 presidential-election estimates from a highly unrepresentative Xbox sample. In a hybrid design combining probability and nonprobability samples, \textcite{dever_combining_2018} found that restricting propensity-score adjustment of a nonprobability sample to the region of common support reduced bias for some estimates, though not all.

But other recent studies reach the opposite conclusion. \textcite{dutwin_apples_2017} found that propensity weighting reduced bias in an opt-in panel but was less effective than simple raking, and adjusted nonprobability samples remained more biased than low-response-rate probability samples. \textcite{yeager_comparing_2011} showed that  demographic weighting improved the accuracy of some non-probability estimates but did not do so for most. Likewise, \textcite{macinnis_accuracy_2018} reported that demographic post-stratification consistently benefited probability samples but rarely improved nonprobability samples. Taken together, this evidence is not reassuring that conventional post-stratification and other post-survey adjustments can solve the inaccuracy of non-probability samples. 

Sample size is another proposed remedy for the error in nonprobability samples, and the ability to obtain large samples quickly and cheaply has repeatedly been presented as an advantage of online nonprobability samples \parencites[12]{baker_report_2013}[24]{yougov_annual_2012}. In probability samples, increasing $n$ reduces the standard error at a well-defined $1/\sqrt{n}$ rate \parencite{cochran_sampling_1977}, with diminishing returns after a few hundred respondents. In nonprobability samples, a larger $n$ does not necessarily reduce selection bias, coverage bias, bogus-response bias, or adjustment-model error \parencite{elliott_inference_2017, mercer_theory_2017}. \textcite{meng_statistical_2018} formalized this concern as the ``big data paradox.'' Using the 2016 CCES, he estimated that, across all eligible voters, the correlation between being included in the data and supporting Trump (Meng's ``data defect correlation'') was about $-0.005$, and showed that at that level, a sample of one percent of eligible voters, about 2.3 million respondents, would have the same mean squared error as a simple random sample of roughly 400. Error that is sampling noise should shrink as respondent count grows; error due to systematic sample bias should persist regardless of how many respondents are added. Thus, the limited evidence currently available does not support the claim that very large sample sizes improve the accuracy of nonprobability samples. 

A handful of studies have examined error in sample matched surveys, specifically, and have found that matching alone provided little or no consistent advantage over conventional post-stratification or raking \parencite{brick_weighting_2015, buskirk_probability_2015, mercer_for_2018}. Evidence on specialized populations similarly indicates that matching can improve sample comparability without necessarily eliminating selection bias \parencite{disogra_matching_2015}.

But not all nonprobability sample survey designs are equivalent. YouGov emerged on the scene using a different method than those used by other vendors. They drew early attention after outperforming traditional telephone pollsters in several UK parliamentary elections, including the 2001 and 2005 contests \parencite{coleman_online_2001, electoralcalculus_2005}. Following its 2007 acquisition of Polimetrix, YouGov entered the US market and adopted Polimetrix’s sample matching methodology. Ahead of the 2008 presidential election, the firm emphasized its ability to deliver rapid, low-cost survey data while maximizing accuracy, positioning sample matching as a potential solution to the limitations of opt-in panels \parencite{crampton_about_2007, bradshaw_yougov_2008}. 

This accuracy claim received notable support in a 2016 Pew Research Center evaluation. In a head-to-head comparison of eight nonprobability vendors and a probability sample fielded with the same questionnaire, YouGov outperformed all other opt-in samples and also exceeded the accuracy of Pew’s probability sample \parencite{kennedy_evaluating_2016}. YouGov exhibited the lowest measurement error across 20 benchmark variables and performed especially well among Black and Hispanic respondents, populations for which survey error is typically higher \parencite{hopkins_getting_2024}. Earlier studies likewise found YouGov samples to produce accurate estimates \parencite{ansolabehere_cooperative_2013, ansolabehere_does_2014, rivers_inference_2009, vavreck_2006_2008}. Therefore, perhaps sample matching represents a substantive advance in nonprobability sample survey design.

However, these past evaluations did not assess the independent impacts of matching and weighting. Each study examined data where both matching and weighting had been done. But since 2014, the CES has incorporated US Senate and gubernatorial election outcomes in its weights, mechanically minimizing error for the races used to demonstrate accuracy. Another limitation of past studies is reliance on a single survey year, which can mislead given the higher variance of nonprobability samples \parencite{macinnis_accuracy_2018}. These limitations motivate a broader, systematic evaluation of CES accuracy and whether conventional demographic post-stratification improves accuracy in nonprobability matched samples independent of YouGov's weighting choices. The investigation described in this paper presents the most comprehensive such assessment to date, using all available election-year surveys administered by YouGov for the CES between 2006 and 2022.

% Subsequent work has evaluated YouGov using election outcomes. \textcite{ansolabehere_cooperative_2013} compared CES vote intentions to state-level election returns from 2006 to 2010, reporting RMSEs between 3.4\% and 4.0 percent across statewide offices. Similarly, \textcite{rivers_inference_2009} found that a CES subsample from 2008 outperformed contemporaneous probability and nonprobability surveys in estimating state-level presidential vote shares.

% (excluded pre jon edit)
% \textcite{rivers_inference_2009} compared a 3,000-respondent CES subsample for the 2008 U.S. presidential election to the ANES Internet Panel, the USA Today–Gallup final pre-election poll, and a Pew Research Center poll, finding that CES had the lowest standardized RMSE (0.06) on state-level estimated party vote share, compared to the ANES panel (0.10), Gallup (0.10), and Pew (0.07). 

% YouGov’s accuracy has not been uniformly high. A 2023 YouGov poll reported unusually high rates of Holocaust denial and opposition to legal abortion among young adults, estimates that diverged sharply from those obtained using probability samples. A replication using a probability-based design attributed these discrepancies to a high prevalence of inattentive or “bogus” respondents in YouGov’s panel \parencite{mercer_online_2024}.

 

% In the early 2000s, prior to its acquisition of Polimetrix, YouGov built its reputation by conducting election polls in the United Kingdom. Twyman (2008) compares the average error of YouGov to other established agencies in predicting the number of seats allocated to each major party across various parliamentary elections. In the 2004 European Parliament election and the 2005 British general election, YouGov performs equivalently to other companies, or marginally better by just about half a percentage point in average error. However, in the 2007 Scottish Parliament election, YouGov's results were slightly less accurate than two other polling agencies by 0.8 and 1.5 percentage points, respectively. Until YouGov, opt-in web surveys were viewed with much skepticism by pollsters and academics. Robert Worcester even called open-access internet surveys "voodoo" polls in \textit{The Independent} in 1985. But these results were viewed favorably, especially given the substantially lower per-respondent cost that YouGov was able to offer. 

\section*{The Cooperative Election Study}
The CES documentation describes its survey samples as representative of US citizens ages 18 and older. The questionnaires have measured an array of variables, including political opinions, demographics, and candidate choices in federal and state elections, including for President, US Senate and House, Governor, Attorney General, Secretary of State, State Senator, and State Representative. Table~\ref{tab:states_respondents_by_office} reports, for each office and year, the number of states with CES respondents and the average number of respondents per state.

\input{output/states_respondents_by_office}

\subsection*{CES Sample Matching and Weighting Procedure}
Sample matching selects respondents from a large opt-in reservoir to resemble a randomly drawn target sample on specified characteristics. For each person in the target sample, the most similar available individual is selected. A sufficiently large, diverse reservoir can approximate the target's joint distribution on those characteristics, which \textcite{rivers_sample_2006} has argued yields a sample that can be treated as random.

Before data collection, YouGov constructed a synthetic sampling frame from the American Community Survey (ACS) and drew a stratified target sample from it. It then invited members of its own opt-in panel and of other companies' panels to complete the CES questionnaire. After data collection, YouGov selected from the pool of completed pre-election interviews the cases that most closely resembled the target sample, a group termed the \textit{matched sample}. The matched sample was then weighted in two steps: first to the sampling frame, by propensity-score weighting (2008--2012) or entropy balancing (2014--2022), and then by post-stratification. The 2006 wave, conducted by Polimetrix, followed a different procedure (Supplementary Material section E). Table~\ref{tab:ces_terms} defines the terms used below and table~\ref{tab:ces_stages} reports the variables used at each stage in each year, as documented in each year's guide (see Supplementary Material section~F for a discussion of documentation errors and inconsistencies).

\input{output/ces_stage_rows}

% In the case of the CES, the matching function relied on overlapping auxiliary information in both the sampling frame and the reservoir, generally demographic variables and sometimes political attitudes and religiosity. Panelists from the reservoir were invited to complete the survey. From the set of respondents who completed the questionnaire, the matching function identified those whose characteristics most closely aligned with the target sample. 

\begin{table}[!htbp]
\centering
\footnotesize
\renewcommand{\arraystretch}{1.35}
\makebox[\textwidth][c]{%
\begin{threeparttable}
\caption{CES sample matching terms.}
\label{tab:ces_terms}
\begin{tabular}{p{4cm} p{10cm}}
\toprule
Term & Definition \\
\midrule
Population of Interest &
All US citizens ages 18 and older. The population about which the CES seeks to draw conclusions. \\[8pt]

Sampling Frame &
A synthetic file of US citizens that YouGov constructed from an earlier year's ACS, with political variables (and, in some years, religious variables) attached from other surveys or modeled. The target sample is drawn from the sampling frame, and the matched sample is weighted to it. \\[8pt]

Target Sample &
A stratified random sample drawn from the sampling frame. Serves as the probability sample that the matched sample attempts to replicate. \\[8pt]

Panel Reservoir &
The members of YouGov's opt-in panel and of the other companies' opt-in panels from which invitees were drawn. \\[8pt]

Invited \newline Respondents &
Panelists selected for invitation to complete the pre-election questionnaire, chosen using a cross-classification on demographic variables.  \\[8pt]

Participating \newline Respondents &
People who completed the pre-election questionnaire (or both the pre-election and post-election questionnaires).  \\[8pt]

Matched Sample &
The participating respondents whom YouGov selected from the pool of completed interviews as most closely resembling the target sample.  \\[8pt]

Frame-Weighted \newline Sample &
The matched sample after weighting to the sampling frame, by propensity score weighting (2008--2012) or entropy balancing (2014--2022). \\[8pt]

Final Weighted \newline Sample &
The frame-weighted sample after post-stratification to demographic and political margins, trimming, and normalization. Its weights are the ones released with the CES. \\[8pt]
\bottomrule
\end{tabular}
\end{threeparttable}%
}
\end{table}

\subsubsection*{The Sampling Frame and Target Sample}

The sampling frame was a stratified sample of citizen records drawn from an earlier year's ACS, with records selected within strata by weighted sampling with replacement. The variables defining these strata are not documented. YouGov then attached political variables, and in some years religious variables, to each frame record. In most years, the guides describe this as linking each frame record to a similar respondent in another survey, such as the Current Population Survey (CPS) or the Pew Religious Landscape Survey, and taking that respondent's values. The sources differed by year (Supplementary Material section E), and the linking process for the CPS data was described as using a weighted distance metric, but its variables and weights were not disclosed.

The target sample was then drawn from the frame by stratification on the variables shown in table~\ref{tab:ces_stages}, with simple random sampling within strata. At this point, YouGov had collected no data; it had only defined the target sample to which CES respondents would later be matched.

\subsubsection*{Panels and Invitations}
Respondents were drawn from opt-in online panels. YouGov has recruited its own panel through advertising and website partnerships, collecting demographics at registration and offering points redeemable for rewards \parencite{yougov_methodology}. From 2008 onward, members of panels run by other companies were also invited (table~\ref{tab:externalproviders}). Invitations were generally allocated across a cross-classification of demographic variables, geography, and sample source, but documentation does not describe how the pool of panelists from which invitees were drawn was determined (Appendix~A, ``Methods used to generate and recruit the sample'').\footnote{An inquiry about the selection process has been made but has not yet been answered.}

\begin{table}[!htbp]
  \centering
  \footnotesize
  \makebox[\textwidth][c]{%
  \begin{threeparttable}
  \caption{Non-YouGov respondent providers.}
  \label{tab:externalproviders}
    \begin{tabular}{l>{\raggedright\arraybackslash}p{10cm}} % Adjust the width as needed
      \toprule
      Year & External Providers \\ 
      \midrule
      2006 & None \\ 
      2008 & E-Rewards, Western Wats \\ 
      2010 & E-Rewards, Western Wats \\ 
      2012 & E-Rewards, Western Wats \\ 
      2014 & MyPoints, Research Now, Survey Sampling International \\ 
      2016 & MyPoints, Research Now, Survey Sampling International, Global Market Insights \\ 
      2018 & Dynata, Prodege, Critical Mix \\ 
      2020 & Dynata, Prodege, Critical Mix \\ 
      2022 & Dynata, Prodege, Disqo, Generation Lab  \\ 
      \bottomrule
    \end{tabular}
  \end{threeparttable}%
  }
\end{table}

\subsubsection*{Data Collection}
In every federal election year since 2006, invited panelists completed a questionnaire before the election, and all who completed it were invited to complete another questionnaire after the election. Field dates appear in table~\ref{tab:sample_description}. Completed pre-election interviews were subjected to quality controls, which the guides do not describe, before entering the pool from which cases were selected.

Table~\ref{tab:sample_description} reports what the CES guides label AAPOR Response Rate 1, which we call a completion rate following AAPOR's standard definitions \parencite[78]{aapor2023}, because the population of interest for an opt-in panel is not well defined. Completion rates ranged from \CompletionRateMin{} percent in \CompletionRateMinYear{} to \CompletionRateMax{} percent in \CompletionRateMaxYear{}, with no monotonic trend, partly because the reporting base changed between years (Appendix~A).

\input{output/sample_description}

\subsubsection*{Selecting Cases from the Pool of Completed Interviews}
The pool of completed pre-election questionnaires exceeded the desired sample size in every year. From the pool, YouGov selected the $N$ respondents who most closely resembled the $N$ cases in the target sample \parencite{rivers_sampling_2007}; the CES guides call this selected set the matched sample. Similarity was measured with a weighted distance metric rather than by exact or propensity score matching \parencite{stuart_matching_2010}. The distance between target case $t$ and respondent $i$ was
\[
D(t,i) = \sum_{k=1}^{K} w_{k} \, d_{k}(t_{k}, i_{k}),
\]
where $d_{k}$ measures dissimilarity on variable $k$ and $w_{k}$ sets that variable's importance, with larger weights demanding closer agreement. For unordered categorical variables, YouGov specified distance matrices between pairs of categories.

Selection brings the joint distribution of these variables in the matched sample close to that of the target sample, but it does not constrain any margin explicitly. Remaining imbalance is addressed by the two weighting steps.

\subsubsection*{Weighting, Step 1: Weighting to the Frame}
The first weighting step adjusted the matched sample toward the sampling frame. From 2008 to 2012, this was done by propensity score weighting. A logistic regression predicted membership in the frame, as against the matched sample, from age, education, gender, and turnout in the previous election; for frame records, turnout was the CPS-reported value attached to the synthetic record. Decile boundaries were set from the distribution of estimated scores in the frame, and matched respondents were weighted so that each decile held the same share of the weighted sample as of the frame.
 
From 2014 through 2022, the matched sample was weighted to the frame by entropy balancing (\cite{hainmueller_entropy_2012}; formal description in Supplementary Material section B). This method assigns each matched respondent a weight such that the weighted sample reproduces specified moments of the frame while departing as little as possible from uniform weights. The balancing conditions always included demographic variables and their interactions, and in some years also political variables (table~\ref{tab:ces_stages}).

\subsubsection*{Weighting, Step 2: Post-Stratification}
The frame-weighted sample was then post-stratified on demographic variables and, in later waves, voter registration, born-again status, and prior presidential vote (table~\ref{tab:ces_stages}). From 2014 onward, weights for the common content were additionally adjusted to the state-level outcomes of that year's gubernatorial and US Senate elections, and in 2020 the post-stratification variables also included that year's presidential vote. Where these outcomes also serve as benchmarks, their agreement with weighted CES estimates partly reflects the weighting itself, a point we return to in the results. Finally, weights were trimmed and normalized to the sample size (Appendix~A).

\section{Benchmarks}
\subsection{Demographics}
Population benchmarks were drawn primarily from Current Population Survey (CPS) data provided by IPUMS \parencite{flood_ipums_2024}. The CPS is a monthly probability sample survey conducted by the US Census Bureau. Annual household response rates ranged from \CpsRateMin{} percent (\CpsRateMinYear{}) to \CpsRateMax{} percent (\CpsRateMaxYear{}) (table~\ref{tab:cps_sample_description}) \parencite{bls_cps_response_rates}, and the data are weighted using census-based population controls (Appendix~A).

% Population targets were primarily drawn from the CPS \parencite{flood_ipums_2024}, generally considered one of the most reliable sources of population statistics \parencite{hur_coding_2013}. The CPS is a probability monthly survey conducted by the U.S. Census Bureau. Each month, approximately 60,000 households are interviewed. CPS data are weighted using a multistage stratified sampling design that adjusts for undercoverage, nonresponse, and demographic characteristics. The weighting process incorporates population controls that are projected from the most recent decennial census and are updated annually. The survey employs a rotation scheme where each household is interviewed for four consecutive months, then dropped from the sample for eight months before potentially being re-interviewed. The interviewing process starts in person, and can be transitioned to telephone after a few months of interviewing. The CPS uses proxy reporting, where one person speaks on behalf of the whole household. 

% Although the response rate has slowly declined over the past two decades, it still remains high at over 70\% (see Table \ref{tab:cesandcps} for annual sample sizes and response rates).


\subsection*{Election Returns and Turnout}

State and federal election returns and voter turnout data were compiled for all 50 states and the District of Columbia from 2006 to 2022 (sources are listed in Appendix~A).

% \indent Comprehensive election returns at both the state and federal levels were compiled from a variety of reputable sources for all 50 states and the District of Columbia for the years 2006 to 2022. Returns for Presidential elections were obtained from open source datasets compiled by the \textcite{mit_election_data_and_science_lab_us_2017-1}. Returns for the U.S. Senate and House of Representatives were sourced from the Federal Election Commission.\footnote{Available at\\ \href{https://www.fec.gov/introduction-campaign-finance/election-results-and-voting-information/}{https://www.fec.gov/introduction-campaign-finance/election-results-and-voting-information/}} Gubernatorial returns were acquired from CQ Press.\footnote{Available with subscription at \href{https://library.cqpress.com/elections/advsearch/election-results.php}{https://library.cqpress.com/elections/advsearch/election-results.php}} Returns for state senate, state representative, secretary of state, and attorney general elections were purchased from Ballotpedia. Voter turnout was compiled by the University of Florida Election Lab \parencite{uf_electionlab_turnout}. Any remaining election results were collected manually from official state websites or the by calling the relevant secretary of state offices.


\section{Methods}
\subsection*{Measuring Error}
Error is the CES-estimated share minus the external benchmark share for the benchmark's modal category or, for vote choice, its winning candidate. The unit of analysis is the state, except for US House contests, which are analyzed by district. Each estimate is compared to the certified share in official election returns. Prior CES evaluations measured vote-choice accuracy using Democratic vote share. The CES asked which \textit{party} respondents voted for in earlier years and which \textit{candidate} in later years, so responses could often not be mapped cleanly to returns (e.g., multiple candidates per party or candidates on multiple party lines). For that reason, we instead use candidate-level comparisons wherever available and party-level comparisons otherwise (Supplementary Material section C lists some problem cases and the years each office had candidate-level data).

Below, we report the root mean squared error (RMSE):

\[
\text{RMSE}_{v,w} = 
\sqrt{\frac{1}{N} \sum_{i=1}^{N} 
\left( \widehat{P}_{\text{CES},\, i,\, w} - P_{\text{true},\, i} \right)^{2}}.
\]
 
Here, $v$ indexes variables, $w$ weighting schemes, and $i$ states or US House districts; $\widehat{P}_{\mathrm{CES},i,w}$ and $P_{\mathrm{true},i}$ are the CES and benchmark shares. $N$ is the number of comparisons. Table~\ref{tab:observation_counts} reports counts by variable and year.

\input{output/observation_counts.tex}
RMSE gives greater weight to large errors, making it a useful summary measure of survey accuracy: surveys that are consistently close to benchmark values across states have lower RMSE than those with similar average accuracy but occasional large errors. Following prior evaluations \parencite{macinnis_accuracy_2018, yeager_comparing_2011}, we distinguish \textit{primary} variables used in either CES weighting step from \textit{secondary} variables used in neither, assigning status by year (table~\ref{tab:ces_stages}). Senate and gubernatorial outcomes from 2014, and presidential vote in 2020, were themselves weighting targets, so estimates for these offices in those years are not independent tests of accuracy.

\subsection{Weighting Scenarios}
Comparing CES accuracy across weighting scenarios allows us to distinguish the error remaining after sample matching from the incremental effect of post-survey adjustment. Because sample matching is intended to approximate a random sample, error in the unweighted CES reflects the performance of matching alone, while differences between weighted and unweighted estimates capture the incremental accuracy provided by weighting.

CES weights reflect design choices specific to YouGov, including the incorporation of known US Senate and gubernatorial outcomes. To assess whether conventional demographic post-stratification improves accuracy in nonprobability matched samples independent of these choices, we evaluated accuracy under an alternative weighting scheme applied consistently across years. 

To this end, we generated an additional set of weights that minimize divergence between the CES and demographic benchmarks from the CPS. Only CES variables with strictly comparable CPS counterparts (identical or fully collapsible categories measuring the same construct) were eligible. Race was excluded because CPS categories could not be exhaustively mapped to CES categories. Although 29 CES variables had conceptual CPS counterparts across years, only \AnesrakeUsedRange{} per year were used for weight generation (table~\ref{tab:vars_used_by_anesrake}). 

\input{output/vars_used_by_anesrake}

Weights were raked to the CPS benchmarks following \textcite{debell_computing_2009} using anesrake \parencite{pasek_anesrake_2018}; the tolerance, weight cap, and convergence criteria are described in Appendix~A.

\section{Results}\label{sec:results}
We document the error from matching alone, how much the CES weights reduce it and whether that depends on how a variable entered the design (through matching, through weighting, or through neither), and then examine changes over time, the distribution of errors, sample size, contest competitiveness, and conventionally computed post-stratification weights.

\subsection{Error with Matching Only}
\input{output/rmse_matching_only}
Table~\ref{tab:rmse_matching_only} reports CES accuracy across all available measures using YouGov's matched sample, but \textit{not} the CES weights. Pooling RMSE estimates across the nine survey waves yields \NObsTotal{} individual error observations. The average RMSE across waves was \MatchYearlyMean{} percentage points, ranging from \MatchYearlyMin{} in \MatchYearlyMinYear{} to \MatchYearlyMax{} in \MatchYearlyMaxYear{}.

Demographic error was stable across waves; the class average had a standard deviation of \MatchDemoSD{} point across waves. Turnout error varied the most (SD = \MatchTurnoutSD{}). Candidate-choice error was steady through 2018, between \MatchCandPreMin{} and \MatchCandPreMax{} points, and then rose to \MatchCandTwentyTwenty{} in 2020 and \MatchCandTwentyTwentyTwo{} in 2022; the CES weights largely removed the 2020 increase (see below).


\subsection{Error with Matching and CES Weights}
Table~\ref{tab:rmse_ces_weighted_all} reports error under the weights YouGov provides, which reflects the error users of the CES would typically encounter.
 
\input{output/rmse_ces_weighted_all.tex}

The weights reduced average error from \MatchYearlyMean{} to \WtdYearlyMean{} percentage points. Substantial error remains in the version of the survey most researchers use. Employment status was the least accurate demographic measure (\WtdEmploymentMean{}). The largest improvements were in candidate choice, whose average fell from \MatchCandMean{} to \WtdCandMean{}; in 2020, it fell from \MatchCandTwentyTwenty{} to \WtdCandTwentyTwenty{}.

Because these averages mix variables inside and outside the weights, the remaining tables classify each variable in each year by whether it was used in matching, in either weighting step, or in neither (2008--2022, because the 2006 documentation does not permit this).

\subsection{Matching Status, CES Weights, and Error}
\input{output/rmse_by_matching_status}
Table~\ref{tab:rmse_by_matching_status} groups variables by whether they were used in matching. Before weighting, variables used in matching had lower error than those that were not (\MatchUsedBefore{} versus \MatchUnusedBefore{}). Matching did not ensure accuracy: employment status was a matching variable yet had an RMSE of \MatchUsedEmploymentBefore{}, and the weights made it slightly worse. The CES weights reduced error by similar amounts in both groups, suggesting that inclusion in matching was less consequential than inclusion in weighting. 

\subsection{Weighting Status, CES Weights, and Error}
\input{output/rmse_by_weighting_status}
By design, weighting should reduce error for the variables used to construct the weights; the more informative test is whether this extends to other variables, as it generally would in a probability sample. Table~\ref{tab:rmse_by_weighting_status} groups variables by whether they were used in either weighting step.
 
Weighting reduced RMSE by \WtUsedReduction{} percentage points for variables used in the weights, compared with \WtUnusedReduction{} points for variables that were not, a gain roughly \WtReductionRatio{} times as large. All \WtUsedRows{} rows for weighting variables improved. Among the \WtUnusedRows{} rows for non-weighting variables, \WtUnusedImproved{} improved and \WtUnusedWorse{} got worse: family income, employment status, union membership, and US House vote.
 
Variables whose status in the CES weights changed across waves show the same pattern. Hispanic origin, turnout, president, US Senate, and governor each improved more in weighting years (those in which the variable was used in the CES weights) than in non-weighting years; each appears in both panels of table~\ref{tab:rmse_by_weighting_status}. For governor and Senate from 2014, and for president in 2020, some of the gain in weighting years is mechanical, because the weights were post-stratified to the benchmark election outcomes. 

\subsection{Error for Variables Used in Neither Matching nor Weighting}

\input{output/rmse_neither_matching_nor_weighting}
The cleanest test of whether the CES design improves accuracy beyond its own targets comes from variables that entered neither matching nor weighting in a given year (table~\ref{tab:rmse_neither_matching_nor_weighting}). For these, weighting reduced average error by \NeitherReduction{} points, from \NeitherBefore{} to \NeitherAfter{}.
 
The CES weights thus provided some benefit beyond the variables used to construct them, but it was smaller and less consistent, and stronger for candidate choice than for demographics. In table~\ref{tab:rmse_neither_matching_nor_weighting}, which pools each office's error across years, every office except the US House improved. Demographic variables outside the design mostly did not improve (family income, employment status, and union membership got worse; Hispanic origin was essentially unchanged), except veteran status and residence duration. Turnout improved but remained in error by \NeitherTurnoutAfter{}. Offices included in the weights generally had lower error and larger apparent reductions from weighting (figure~\ref{fig:rmse_reduction_by_office}).

\begin{figure}
    \centering
    \includegraphics[width=\linewidth]{output/rmse_reduction_by_office.png}
    \caption[RMSE reduction from CES weighting by office]{RMSE reduction from CES weighting by office. Bars show matching-only RMSE and the RMSE remaining after CES weighting, pooled over all error observations from the 2008--2022 waves, as in tables~\ref{tab:rmse_by_weighting_status} and~\ref{tab:rmse_neither_matching_nor_weighting}. An office appears among primary offices in the years it was a CES weighting target and among secondary offices otherwise. The 2006 wave is excluded because matching and weighting classifications are unavailable.}
    \par\smallskip
    \noindent\begin{minipage}{\linewidth}\footnotesize\textit{Alt text:} Bar chart of error before and after weighting for eight offices. Weighting cut error by up to \WtUsedPresReduction{} percentage points when an office was a weighting target, by \NeitherOfficeReductionMin{} to \NeitherOfficeReductionMax{} points otherwise, and slightly raised United States House error.\end{minipage}
    \label{fig:rmse_reduction_by_office}
\end{figure}
 
% The preceding results combine primary and secondary variables. Post-stratification should, by design, reduce error for primary variables (since they are used to construct the weights in the first place). The more informative test is whether those gains extend to secondary variables in the matched sample, as they generally would in a probability sample. 

% \begin{figure}
%     \centering
%     \includegraphics[width=\linewidth]{output/primary_secondary_matching_vs_weighting.png}
%     \caption{Effect of CES Post-Stratification on Primary and Secondary Variable Error}
%     \caption*{\footnotesize
%     Bars show mean RMSE for matching-only and matching + post-stratification estimates by year. Panel A includes Primary variables used to construct the CES weights; Panel B includes Secondary variables not used to construct the weights. Labels report $\Delta$ RMSE, calculated as post-stratified RMSE minus matching-only RMSE, so negative values indicate improved accuracy after weighting. Only variables consistently classified as Primary or Secondary and present in every CES wave from 2008 through 2022 are included.}
%     \label{fig:primary_secondary_weighting_comparison}
% \end{figure}

% \input{output/primary_secondary_weighting_comparison}


% Illustrated in Figure~\ref{fig:primary_secondary_weighting_comparison}, post-stratification reduced average primary-variable error from 6.9 to 3.7 percentage points on average, a 46 percent reduction, and it did so in every year. Average error on secondary variables fell from 7.7 to 7.0 points, a 10 percent reduction, and in six of eight years the improvement was less than a single percentage point. Averaged across years, weighting improved primary variables by 3.2 points and secondary variables by 0.8 — a fourfold difference (Table~\ref{tab:primary_secondary_weighting_comparison}). The benefits of post-stratification  did not extend to variables outside the weighting scheme, consistent with other types of nonprobability sampling. Before post-stratification, primary and secondary variables were estimated with comparable accuracy; afterward, error for the variables used to construct the weights was roughly half the error for those that were not — 3.7 percent versus 7.0.

% This gap does not appear to be explained by primary variables having started further from the benchmark. Unweighted error was in fact slightly higher for secondary variables (7.7) than for primary ones (6.9), so the set that improved less began, if anything, with more error to correct.

% Nor does it appear to reflect variable composition. The consistently primary variables are three demographics — age group, education, and sex — while the secondary set mixes demographic and candidate-choice measures. We therefore repeated the comparison two ways. Within demographic variables alone, which begin at nearly identical unweighted error (6.9 and 7.0 points), post-stratification reduced primary error by 3.2 points and secondary error by 0.2. For candidate choice, we can look at governor and U.S. Senate, which swapped status: they were secondary from 2008 through 2012 and primary thereafter. Post-stratification reduced their error by 3.7 points in the years they were primary and by 1.8 points in the years they were secondary. Both comparisons are reported in Appendix Table~\ref{tab:demographic_candidate_choice_weighting_comparison_vertical}.
%% PREVIOUS VERSION AS OF 9/5/26 ENDS HERE

% Variables included in CES post-stratification (primary variables) exhibited lower error than variables that were not (secondary variables; see Table~\ref{tab:finalpoststratvars}). As shown in Figure~\ref{fig:primary_vs_secondary_consistent_excl_2006} and Table~\ref{tab:primary_secondary_rmse_by_year}, average RMSE for primary variables was 3.7 percent, compared to 7.0 percent for secondary variables, yielding an average within-year gap of 3.2 percentage points (Figure~\ref{fig:primary_vs_secondary_delta_by_year}). These calculations include only those variables that were either primary or secondary in each of the nine survey years examined and excludes any variable that changed primary/secondary status. These calculations include only variables that held the same status—primary or secondary—in all nine survey years examined. Candidate choice for governor and U.S. Senate, along with voting turnout, is therefore excluded, since each changed status in 2014.

% \footnote{Because the errors for turnout were so much larger than all the other errors, and because turnout was a primary variable in some CES years and a secondary variable in other years, turnout was excluded from these calculations.}

% \begin{figure}
%     \centering
%     \includegraphics[width=\linewidth]{output/primary_vs_secondary_consistent_excl_2006.png}
%     \caption{Primary vs. Secondary Variable RMSE}\label{fig:primary_vs_secondary_consistent_excl_2006}
%      \caption*{\footnotesize Primary variables (3): Age Group, Education, Sex. Secondary variables (8): Family Income, Employment Status, State Representative, State Senator, U.S. House, Union Membership, Veteran Status, and Voting Method. Only variables consistently classified as Primary or Secondary across all years and present in every CES wave are included. 2006 is excluded because the CES does not report post-stratification variables for that year.   
%     }
% \end{figure}

% \input{output/primary_secondary_rmse_by_year.tex}

% \begin{figure}
%     \centering
%     \includegraphics[width=\linewidth]{output/primary_vs_secondary_delta_by_year.png}
%     \caption{Difference between Primary and Secondary by Year} 
%     \caption*{\footnotesize Positive values indicate Secondary variables had higher RMSE than Primary variables.}
%     \label{fig:primary_vs_secondary_delta_by_year}
% \end{figure}

\subsection{Distribution of Errors}
Prior work has shown that a hallmark of nonprobability samples is wide variation across measures in the amount of error observed. It has further shown that post-stratification, which corrects known imbalances in probability samples, produces smaller and less consistent gains in nonprobability ones \parencite{yeager_comparing_2011, macinnis_accuracy_2018, chang_national_2009, lavrakas_comparing_2022, walker_foundations_2009, mercer_for_2018}.

\begin{figure}
    \centering
    \includegraphics[width=\linewidth]{output/error_distribution_matching_vs_weighting.png}
    \caption[Distribution of CES absolute error before and after weighting]{Distribution of CES absolute error before and after weighting. Each histogram shows the share of all errors at each level of percentage-point absolute error, pooling \NObsTotal{} observations across \NVariables{} variables and nine CES waves. Panel A reports matching-only estimates; Panel B reports estimates using YouGov's weights. Weighting concentrated more estimates near their benchmarks but left the overall dispersion of errors essentially unchanged (SD = \ErrSdMatching{} and \ErrSdWeighted{}, respectively).}
    \par\smallskip
    \noindent\begin{minipage}{\linewidth}\footnotesize\textit{Alt text:} Two histograms of absolute error before and after weighting. After weighting, more estimates sit near zero, but both distributions keep a long tail of errors above 20 percentage points.\end{minipage}
    \label{fig:error_distribution_matching_vs_weighting}
\end{figure}

Figure~\ref{fig:error_distribution_matching_vs_weighting} shows both patterns across the 2006–2022 CES waves. Under matching alone, absolute error was not concentrated near zero: \ErrWithinOnePtMatching{} of estimates were within one percentage point of their benchmark, \ErrWithinThreePtMatching{} were within three points, and \ErrBeyondFivePtMatching{} missed their benchmark by more than five points. The CES weights shifted the distribution toward the origin, increasing the share of estimates within three points from \ErrWithinThreePtMatching{} to \ErrWithinThreePtWeighted{}. Mean absolute error declined from \ErrMeanMatching{} to \ErrMeanWeighted{} percentage points, and the median declined from \ErrMedianMatching{} to \ErrMedianWeighted{} points. The standard deviation of absolute error, however, was nearly unchanged at \ErrSdMatching{} points under matching alone and \ErrSdWeighted{} points after weighting, suggesting that weighting improved the typical estimate without producing a comparable contraction in the extreme tail.

The effect of applying the CES weights was smaller for larger errors. The proportion of errors exceeding 15 points declined only from \ErrBeyondFifteenMatching{} to \ErrBeyondFifteenWeighted{} after weighting, while the proportion exceeding 20 points changed from \ErrBeyondTwentyMatching{} to \ErrBeyondTwentyWeighted{}. In total, \ErrBeyondTwentyBoth{} of estimates exceeded 20 points under both methods. Of the \NBeyondTwentyMatching{} estimates whose matching-only errors exceeded 20 points, \NMovedBelowTwenty{} moved below that threshold after weighting and \NStayedBeyondTwenty{} remained above it. At the same time, \NCrossedAboveTwenty{} estimates that had previously been below 20 points crossed above the threshold, leaving \NBeyondTwentyWeighted{} weighted estimates above it. The CES weighting procedure therefore changed the composition of the extreme-error tail but produced little change in its overall frequency.

\subsection{Error Across Time}
Advocates of nonprobability samples have sometimes argued that their accuracy has improved as opt-in online panels expanded, recruitment methods matured, and weighting procedures became more sophisticated. The CES offers a natural test of this claim: over the period studied, it revised its weighting scheme repeatedly and drew respondents from multiple panel providers.

We estimated linear trends for primary and secondary variables separately, to see whether error changed differently over time depending on whether a variable was used in weighting (figure~\ref{fig:rmse_trend_primary_secondary}; table~\ref{tab:rmse_trend_primary_secondary}). To keep each group's composition fixed, we used only variables whose status never changed: the three variables that were primary in every documented wave from 2008 through 2022 (age group, sex, and education) and the eight that were secondary in every wave from 2006 through 2022, with the undocumented 2006 wave treated as secondary. For each group, we regressed year-by-variable RMSE on calendar year with variable fixed effects, separately for matching-only and weighted estimates, using HC3 robust standard errors and a permutation test that randomly reassigned year labels.

For the primary variables, matching-only error showed no trend over the period (\TrendPrimMatchBeta{} percentage points per year; SE = \TrendPrimMatchSE{}, p = \TrendPrimMatchP{}), so the matched samples did not become any more representative, over time, on the variables used to construct the weights. Weighting did produce an apparent decline (\TrendPrimWtdBeta{} points per year; SE = \TrendPrimWtdSE{}, p = \TrendPrimWtdP{}; permutation p = \TrendPrimWtdPermP{}), but it is concentrated early in the period. Mean weighted error was \TrendPrimWtdMeanFirst{} in 2008 and \TrendPrimWtdMeanTwentyTen{} in 2010, then settled between \TrendPrimWtdLateMin{} and \TrendPrimWtdLateMax{} in every wave from 2012 onward.

For the secondary variables, the pattern differed. Under matching alone, secondary-variable error rose by \TrendSecMatchBeta{} points per year (SE = \TrendSecMatchSE{}, p = \TrendSecMatchP{}; permutation p = \TrendSecMatchPermP{}), an estimated increase of about \TrendSecMatchFullPeriod{} points from 2006 to 2022. Weighting offset only part of this increase: with weighting, the trend was smaller (\TrendSecWtdBeta{} points per year) and no longer statistically significant (SE = \TrendSecWtdSE{}, p = \TrendSecWtdP{}; 95\% CI: [\TrendSecWtdCILow{}, \TrendSecWtdCIHigh{}]; permutation p = \TrendSecWtdPermP{}), but the point estimate remained positive.

For the variables used to build the weights, matching-only error showed no improvement, and the apparent decline after weighting reflected the weights hitting their own targets more precisely, not a more representative sample underneath. For the variables not used to construct the weights, error from matching alone rose, and weighting offset only part of that increase. Our results suggest that the CES matched samples did not become more accurate over the 17-year period we studied.

\begin{figure}
    \centering
    \includegraphics[width=\linewidth]{output/rmse_trend_primary_secondary.png}
    \caption[RMSE trends over time for primary and secondary variables]{RMSE trends over time for primary and secondary variables. Each point is the mean RMSE across a group's variables in one wave; solid lines are fitted trends from variable fixed-effects regressions of year-by-variable RMSE on calendar year, averaged across variables, and faint lines connect the yearly means. Panel A: primary variables (age group, sex, and education), used in weighting in every documented wave, 2008--2022. Panel B: secondary variables (family income, employment status, union membership, veteran status, voting method, US House, State Senator, and State Representative), never used in weighting, 2006--2022; because the 2006 guide does not document the matching or weighting variables, 2006 is included on the assumption that these variables were secondary. Labels report the slope $b$ in percentage points per year and its HC3 $p$-value; table~\ref{tab:rmse_trend_primary_secondary} reports full results.}
    \par\smallskip
    \noindent\begin{minipage}{\linewidth}\footnotesize\textit{Alt text:} Two scatterplots of error by year with fitted trends. For weighting variables, matching-only error is flat and weighted error falls. For other variables, matching-only error rises from about \TrendSecMatchMeanFirst{} to \TrendSecMatchMeanLast{} percentage points, and weighted error rises less.\end{minipage}
    \label{fig:rmse_trend_primary_secondary}
\end{figure}

\input{output/rmse_trend_primary_secondary_tests.tex}

% \subsubsection{The Relationship Between Error and Sample Size}
% Comparing the sample size gradient respondent-count gradient in the CES to the $1/\sqrt{N}$ pattern from probability sampling indicates how much of the observed benchmark error was plausibly reducible through larger samples, as opposed to error that persisted even after sampling variability should have declined.

% Table~\ref{tab:respondent_count_error_by_variable} reports separate OLS regressions, estimated for each variable, of absolute percentage-point benchmark error on $1/\sqrt{N}$, where $N$ is the number of CES respondents. The unit of observation is a state-year for statewide measures and a district-year for district-level measures; all regressions include state and year fixed effects.

% This specification evaluates whether CES error follows the rate implied by ordinary probability-sample sampling error. Under that benchmark, error should increase approximately linearly with $1/\sqrt{N}$: estimates based on fewer respondents should have larger errors.

% \input{output/respondent_count_error_by_variable.tex}

% The results show that respondent count matters, but unevenly. The clearest relationships appear for candidate-choice variables, especially offices whose estimates are built from small geographic or contest-specific respondent pools. U.S. House and State Representative estimates show statistically significant respondent-count gradients under both specifications, and several other candidate-choice offices have coefficients in the expected direction. Demographic variables, by contrast, show weak, inconsistent, and generally insignificant relationships with respondent count. 

% Figure~\ref{fig:respondent_count_sqrt_fit} visualizes the same question, plotting observed RMSE across sample size bins and overlaying a fitted curve of the form
% \[
%     \text{RMSE}(N) = \alpha + \frac{\beta}{\sqrt{N}}.
% \]
% Here $\alpha$ is the estimated error floor: because $\beta/\sqrt{N} \to 0$ as $N$ grows, the curve approaches $\alpha$ at large respondent counts. The parameter $\beta$ governs the sample-size-scaling component, with larger values implying that RMSE falls more steeply as respondent count increases. The curve is descriptive rather than structural: it asks whether the empirical pattern resembles the conventional sampling-error relationship while allowing error to flatten at a nonzero floor. Accordingly, $\alpha$ should be read as an empirical residual-error floor, not a theoretical lower bound.

% \begin{figure}
%     \centering
%     \includegraphics[width=\linewidth]{output/observed_rmse_vs_inverse_sqrt_n_shape_by_class_excl_turnout_unweighted.png}
%     \caption{Observed and Extrapolated RMSE by Sample Size, Compared to Theoretical Probability Sample}
%     \caption*{\footnotesize
%     Points show observed RMSE across sample-size bins for all variables except turnout. The solid gold curve shows the fitted relationship $\alpha + \beta/\sqrt{N}$ over the observed sample-size range, while the dashed gold segment extrapolates that relationship to $N=25{,}000$. The gray dash-dot curve shows the theoretical $1/\sqrt{N}$ decline expected under probability sampling, calibrated to the same RMSE at $N=50$; the dotted horizontal line marks the fitted nonzero error floor. Sample size is displayed on a logarithmic scale. Results use unweighted CES estimates. Turnout is excluded because its substantially higher error exhibits little relationship with sample size, suggesting that systematic factors predominate and causing the pooled $1/\sqrt{N}$ relationship to collapse.}
%     \label{fig:respondent_count_sqrt_fit}
% \end{figure}


% A larger sample size was associated with lower error, but the relationship was noisy and bounded. RMSE declined as respondent count increased and mirrored the $1/\sqrt{N}$ shape expected under probability sampling, but the relationship flattened well above zero, with a fitted floor of approximately 6.6 percentage points. 

\subsection{Sample Size and Estimation Error}
Because the CES estimates quantities for states and districts whose respondent counts vary within and across waves, the sample sizes underlying its error estimates range from \EstSampleSizeMin{} to \EstSampleSizeMax{} respondents (Supplementary Material section A, figure~\ref{fig:app_sample_size_distribution}). This variation lets us assess whether error in these matched samples declines with sample size as it would in a probability sample.

Figure~\ref{fig:respondent_count_sqrt_fit} provides a descriptive visualization of the relationship between error and sample size for variables other than turnout. It plots observed RMSE across sample-size bins and overlays a curve of the form
\[
    \operatorname{RMSE}(N)
    =
    a+\frac{b}{\sqrt{N}}.
\]
The curve is constrained to begin at the observed RMSE for the first bin, approximately \AnchorRmse{} percentage points at $N=50$. Under this constraint, $a$ is estimated from the binned data and $b$ is implied by the common starting point. Because $b/\sqrt{N}\rightarrow 0$ as $N$ increases, the fitted curve approaches $a$ at large sample sizes, making $a$ a descriptive residual-error floor for these matched, nonprobability samples. 

\begin{figure}[!htbp]
    \centering
    \includegraphics[width=\linewidth]{output/observed_rmse_vs_inverse_sqrt_n_shape_by_class_excl_turnout_unweighted.png}
    \caption[Observed and extrapolated RMSE by sample size relative to the probability-sampling benchmark]{Observed and extrapolated RMSE by sample size relative to the probability-sampling benchmark. Points show observed RMSE across sample-size bins for all variables except turnout. Bins are sample-size quantiles of roughly \FitBinObs{} error observations each. Each point is a single RMSE over all observations in its bin, pooled across waves and variables. The solid curve shows the fitted relationship $a+b/\sqrt{N}$ over the observed sample-size range, while the dashed segment extrapolates that relationship beyond the available data to $N=25{,}000$. The dash-dot curve shows the theoretical $1/\sqrt{N}$ decline expected under probability sampling, calibrated to the same RMSE at $N=50$. The dotted horizontal line marks the fitted residual-error floor. For comparison, a simple random sample estimating a proportion near one half would have an expected error equal to the fitted floor at about \SrsNAtFloor{} respondents. Results use unweighted CES estimates. Turnout is excluded from this visualization because it has substantially higher error.}
    \par\smallskip
    \noindent\begin{minipage}{\linewidth}\footnotesize\textit{Alt text:} Line chart of error against sample size on a logarithmic scale. Observed error falls from \AnchorRmse{} percentage points at 50 respondents and levels off near a fitted floor of \ErrorFloor{} points, while the probability-sampling curve keeps falling toward zero.\end{minipage}
    \label{fig:respondent_count_sqrt_fit}
\end{figure}

Observed RMSE initially declined with sample size, but the fitted relationship flattened at approximately \ErrorFloor{} percentage points, whereas the probability-sampling curve continues toward zero.

\subsection{Error as a Function of Contest Competitiveness}
% The analysis so far pools respondents within state-years. Because U.S. House estimates are constructed at the level of the congressional district rather than the state, they let us check whether the same respondent-count relationship holds at a finer level of aggregation. Figure~\ref{fig:us_house_error_density_combined} plots district-level U.S. House mean absolute error against per-district respondent count, alongside the distribution of respondents across districts. The district-level pattern mirrors the state-level one: error declines as the number of respondents per district rises.

% \begin{figure}
%     \centering
%     \includegraphics[width=\linewidth]{output/us_house_error_density_combined.png}
%     \caption{U.S. House Mean Absolute Error and Respondent Density by District}
%     \label{fig:us_house_error_density_combined}
% \end{figure}

As the winning candidate's benchmark vote share moved farther from 50 percent---that is, as races became less competitive---RMSE generally increased, contrary to the prediction of sampling variance alone, which is largest at an even split. The increase was most pronounced in the least competitive races, including races where a major-party candidate dropped out or only one appeared on the ballot, where the benchmark winning share could approach or exceed 90 percent while the survey-implied split remained much closer to parity. For most offices, RMSE is relatively stable through moderate winning-share cutoffs but rises once high-share races are included, most visibly for Attorney General, Secretary of State, and State Senator (table~\ref{tab:office_rmse_less_competitive_races}; figure~\ref{fig:office_rmse_progressive_race_exclusion}).

\input{output/office_rmse_progressive_race_exclusion.tex}

\begin{figure}
    \centering
    \includegraphics[width=\linewidth]{output/office_rmse_progressive_race_exclusion.png}
    \caption[Sensitivity of office-level RMSE to less competitive races]{Sensitivity of office-level RMSE to less competitive races. Each point recalculates office-level RMSE using only races where the winning candidate's benchmark vote share is at or below the x-axis cutoff. For example, 100 percent includes all contests for that office in the RMSE calculation, whereas 80 percent excludes races where the winning candidate received more than 80 percent of the vote. Moving from left to right progressively includes less competitive races. Error generally increases as less competitive contests are retained, with the sharpest increases occurring after very high winning-share races are included.}
    \par\smallskip
    \noindent\begin{minipage}{\linewidth}\footnotesize\textit{Alt text:} Line chart of error for eight offices as less competitive races are added. Error stays nearly flat for President, United States Senate, Governor, and United States House, and rises most for Attorney General once races won with over 90 percent are included.\end{minipage}
    \label{fig:office_rmse_progressive_race_exclusion}
\end{figure}

This compression toward parity contrasts with prior work documenting overstatement of winners, incumbents, and frontrunners in post-election reports and pre-election polls \parencite{wright1993errors, mattei1998winning, hopkins2009wilder, chen2025electoral}. Possible mechanisms are that respondents reported their party's candidate rather than a crossover vote for a strong winner, less salient contests invited weaker recall, and turnout overreporting added vote choices from nonvoters.

US House races let us examine sample size and competitiveness jointly. Because the unit of inference was the district and many districts contained relatively few CES respondents, House contests combined both sources of error in a way that state-level estimates for other offices did not. Table~\ref{tab:us_house_rmse_grid} presents a two-way RMSE matrix that stratifies US House races by the winner's vote share and by minimum district-level respondent thresholds, showing how competitiveness and district sample size interact.
% \begin{figure}
%     \centering
%     \includegraphics[width=\linewidth]{output/rmse_by_competitiveness_histogram.png}
%     \caption{RMSE Across Different Bins of Race Competitiveness, All Races}
%     \caption*{\footnotesize
%     Competitiveness refers to the distance from the winning candidate receiving 50\% of the vote. Results are matching + CES post-stratification. Includes all races: President(n=204), U.S. Senate (n=303), U.S. House (n=2,888), Governor (n=225), Attorney General (n=175), Secretary of State (n=146), State Rep. (n=397), State Senator (n=367)}  
%     \label{fig:rmse_by_competitiveness_histogram}
% \end{figure}

% CES-weighted accuracy varied based on the competitiveness of the race. The CES was less accurate across all races than when only competitive races were considered. Figure \ref{fig:rmsebyracecomp} compares the accuracy of the CES for all non-weight-generation races to those non-weight-generation races for which the actual Democratic vote share was between 40 and 60 percent, indicating a tighter election. The deltas correspond to the difference between the all-inclusive average and the close-race average. In all years, it was the case that the close-race RMSE was lower than the all-inclusive average. Across years, the RMSE was 2.3\% greater across all races than only in more competitive races. 

% This is perhaps counterintuitive. One might reasonably expect it to be harder to accurately measure closer races because the smaller margin of difference makes outcomes more uncertain and sample variations potentially more impactful. However, non-competitive races are typically less visible to voters. When respondents are asked about vote choices for constitutional offices such as secretary of state or attorney general, they may have been less familiar with the candidates and more likely to report having voted along party lines. In races where candidates ran unopposed, many respondents still reported voting for a party, inducing large errors in these states since around half of the respondents claimed to vote for a party that received no votes.
\subsection{Comparison to Prior Evaluations of Error in the CES}

CES documentation reported accuracy metrics for selected offices and years, though these are narrower in scope and do not distinguish between primary and secondary variables. Table~\ref{tab:ces_comparison} compares the RMSEs reported in the guides with those computed in this study.

\input{output/ces_comparison}

Our RMSEs are generally higher than the reported values. The difference is small for President (\GuideGapPres{} points on average) and negligible for US Senate (\GuideGapSenate{}), but larger for Governor (\GuideGapGov{}), Secretary of State (\GuideGapSoS{}), and Attorney General (\GuideGapAG{}). Averaged across years, our estimates exceed the reported values by \GuideGapOverall{} points.

We cannot determine with certainty what accounts for this difference, but several features of the guides' approach may contribute. The guides measured the Democratic share of the two-party vote, whereas we use the modal candidate's share of all votes cast; they used validated-voter weights from 2018 onward; and their 2006 and 2010 editions included only states with at least 200 respondents. In addition, certified vote shares for the same race differ somewhat across sources.
% \begin{sidewaystable}[p]
% \centering
% \caption{CES-Reported RMSEs vs. RMSEs Calculated in This Study}
% \label{tab:ces_comparison}
% \input{output/ces_comparison.tex}
% \end{sidewaystable}

\subsection{Effect of Conventional Post-Stratification on CES Error}
We then evaluated how closely standard best-practice post-stratification of the nonprobability sample approached the accuracy of YouGov's bespoke weighting scheme, using a single national set of \texttt{anesrake} calibration weights and, for a fair comparison, evaluating accuracy only on variables used in neither YouGov's weighting nor anesrake.

As a validation, RMSE on the calibration variables averaged at most \AnesrakeFitMeanMax{} percentage points across years (table~\ref{tab:national_anesrake_variables_accuracy}), indicating alignment with national benchmarks. This accuracy did not extend to variables excluded from both weighting schemes: State Representative, State Senator, and US House races (figure~\ref{fig:rmse_anesrake_full_by_year}). Averaged across these three offices, anesrake performed worse than both matching alone and matching plus the CES weights, except in 2020 and 2022, when it improved on or roughly equaled matching alone but still underperformed YouGov's weights. By office, anesrake was more accurate than matching alone for both state-legislative offices in 2014, 2020, and 2022 but less accurate for US House in every year (table~\ref{tab:rmse_anesrake_full}). That table also shows that YouGov's weights reduced RMSE by an average of \CesThreeVarReduction{} percentage points across the three variables not used for weighting, whereas anesrake increased it by \AnesThreeVarIncrease{} percentage points, consistent with prior findings that best-practice demographic weighting does not generalize to unconstrained variables in nonprobability sample surveys \parencite{yeager_comparing_2011, macinnis_accuracy_2018}.

\begin{figure}
    \centering
    \includegraphics[width=\linewidth]{output/rmse_anesrake_full_by_year.png}
    \caption[Comparison of weighting techniques across years]{Comparison of weighting techniques across years. Each bar is the mean RMSE across three variables in one wave, computed from the same respondents under three weighting scenarios: matching only, with no weights; matching plus the weights YouGov releases with the CES; and matching plus national calibration weights we constructed with anesrake to CPS demographic benchmarks (table~\ref{tab:vars_used_by_anesrake}). The three variables, US House, State Senator, and State Representative vote choice, are the only ones observed in every wave and used in neither the CES weights nor anesrake, so neither weighting scheme was fit to them. For 2006, whose guide does not document the CES weighting variables, we assume they were not used.}
    \par\smallskip
    \noindent\begin{minipage}{\linewidth}\footnotesize\textit{Alt text:} Grouped bar chart of mean error for three down-ballot offices by year under three weighting approaches. Anesrake weights gave the highest error in \AnesHighestYears{} of \ThreeVarYears{} years, and the Cooperative Election Study weights cut error most in 2020 and 2022.\end{minipage}

    \label{fig:rmse_anesrake_full_by_year}
\end{figure}

\input{output/rmse_anesrake_full_table.tex}


\section{Discussion}\label{sec:discussion}
This study provides the most comprehensive evaluation to date of matched-sample survey accuracy. Across nine waves of the CES and \NObsTotal{} benchmark comparisons, error averaged \MatchYearlyMean{} percentage points under matching alone and \WtdYearlyMean{} with the CES weights. Accuracy did not improve over the 17 years studied. For variables never used in the weights, matching-only error rose by about \TrendSecMatchFullPeriod{} points from 2006 to 2022, while the weighted trend was smaller and not significant. This held despite changes of panel providers and procedural revisions.

A central finding is that the accuracy of the CES depends more on YouGov's weights than on sample matching. The theoretical case for matching is that, if selection into the panel is ignorable given the matching variables, a matched sample can be treated as a random sample \parencite{rivers_sample_2006, rivers_sampling_2007}. In the CES, variables outside matching were estimated with substantial error, which suggests that the matching variables did not fully account for how respondents differed from the population \parencite{mercer_theory_2017}. Matching also did not ensure accuracy for the variables it used: employment status was a matching variable in 2014 and 2016 yet remained among the least accurate measures in the study, and whether a variable was used in matching made little difference to how much the weights improved it. In practice, matching in the CES appears to function less as a substitute for random selection than as a first-stage adjustment whose shortcomings the weights must then correct.

The CES weights, in turn, improved accuracy most for the variables used to construct them. This is consistent with prior work showing that weighting improves an outcome mainly insofar as it is related to the weighting variables \parencite{mercer_for_2018}. The CES weights draw on political information, including presidential vote and, from 2014, statewide election results, which may explain why their benefits reached other vote choices more consistently than demographic characteristics. By contrast, conventional demographic post-stratification with anesrake increased average error across the three vote-choice outcomes observed in every wave and excluded from both weighting schemes, replicating prior evidence that demographic post-stratification does not generalize to unconstrained variables in nonprobability samples \parencite{yeager_comparing_2011, macinnis_accuracy_2018, dutwin_apples_2017}.

These findings align with prior evidence that average error in nonprobability samples is often only modestly worse than in probability samples but that variance is substantially greater, including a nontrivial share of large errors \parencite{yeager_comparing_2011, chang_national_2009, macinnis_accuracy_2018, cornesse_review_2020}. The weights improved the typical estimate but left the dispersion of errors essentially unchanged, with roughly one estimate in ten still off by more than 15 points, so the relevant risk for a researcher analyzing a single state or office is a badly wrong estimate, not the average error. Errors were also largest in the least competitive races, where CES estimates were compressed toward an even split rather than overstating winners as earlier election surveys did \parencite{wright1993errors, mattei1998winning}.

The sample-size results point the same way. In the CES, error declined with respondent count but leveled off at about \ErrorFloor{} percentage points for variables other than turnout, rather than approaching zero as in probability samples \parencite{cochran_sampling_1977}. The fitted curve implies that most of the reducible error is gone early: by roughly 500 respondents, predicted error is within about \FloorGapAtFiveHundred{} points of the floor. This is the pattern expected when the remaining error is bias, which larger samples do not remove \parencite{meng_statistical_2018, bethlehem_selection_2010}. The distinction matters because large samples are what nonprobability designs make cheap \parencite{baker_report_2013}, raising the question of whether a large nonprobability sample is preferable to a smaller probability sample of the same population.

Nonprobability sampling is now a permanent feature of the survey research landscape. YouGov occupies a prominent position in academic and applied research, and its use of sample matching has been presented as a solution to the limitations of opt-in panels \parencite{rivers_sampling_2007, crampton_about_2007}. Our evidence challenges that interpretation. The error remaining after sample matching in the CES is comparable to that reported for other nonprobability sampling methods in prior evaluations. The CES's accuracy instead depends heavily on weights that incorporate outcome-adjacent information, such as prior vote and election results, an approach that is unavailable for many research applications that lack benchmarks. For those applications, our results for variables used in neither matching nor weighting are the more relevant guide to expected accuracy.

This distinction also bears on how CES accuracy has been assessed. Prior evaluations, including those in the CES guides, have largely compared estimates with election outcomes \parencite{ansolabehere_cooperative_2013, ansolabehere_does_2014, rivers_inference_2009, vavreck_2006_2008}. We reproduce lower error estimates in some cases as well, but we show that these conclusions are sensitive to which variables are evaluated and whether they were used in weighting. From 2014 onward, statewide Senate and gubernatorial results were weighting targets, so agreement with them is not independent evidence of sample quality.

Several limitations qualify these findings. Incomplete documentation required year-specific classifications. We exclude 2006 from every comparison that sorts variables by whether they were used in matching or weighting. We include 2006 only in the trend and anesrake analyses, and only for variables that were never used in the CES weights in any documented year, which we assume were also unused in 2006. Few variables were consistently outside the weights, and many were vote choices. Special elections, fusion voting, ballot changes, and benchmark definitions complicate candidate comparisons (Supplementary Material section C). Turnout overstatement also occurs in probability sample surveys \parencite{debell_turnout_2020, hammer_experiments_2014}, so it is not wholly attributable to nonprobability sampling. These limitations reflect real constraints faced by researchers using CES data, but they affect a small share of the total errors observed, making it unlikely that they explain the broader patterns we observe.

The high variance of nonprobability samples implies that single-year evaluations, still common, can mislead; multi-year designs are essential for assessing stability, trends, and the frequency of large errors. Evaluations should report error separately for variables inside and outside the weights, which producers could support by documenting the variables used at each stage of each wave. Understanding why weighting that reproduces national margins fails to improve state- and district-level estimates may shed light on the sources of error. Comparative evaluations across vendors remain rare (the Pew comparison \parencite{kennedy_evaluating_2016} is one of the few that includes YouGov), and more would be valuable.

In sum, our findings challenge the view of sample matching as a solution to the accuracy problems of nonprobability sample surveys. Combined with YouGov's weights, matched samples perform best on the variables the weights target and somewhat better on related political measures, but they remain subject to the constraints that characterize nonprobability designs: bias that larger samples do not remove, a persistent share of large errors, and limited generalizability to outcomes the weights do not target. Researchers using CES data should interpret results with caution, particularly for variables not targeted by the weighting process.

% In probability samples, the standard error of an estimate falls in proportion to $1/\sqrt{n}$ as $n$ increases. In nonprobability samples such as the CES, a larger $n$ still reduces random variability, but it does not necessarily reduce selection bias, coverage bias, bogus-response bias, or adjustment-model error \parencite{elliott_inference_2017, mercer_theory_2017}. This distinction matters because nonprobability samples are precisely the ones for which large $n$ is cheap to obtain, which raises the question of whether a large nonprobability sample is preferable to a smaller probability sample of the same population.

% This study provides the most comprehensive evaluation to date of the accuracy of matched-sample surveys using YouGov data from the Cooperative Election Study. Across nine survey waves spanning 2006–2022, CES estimates exhibited an average RMSE of approximately 8 percent, with errors reaching 10–15 percent in some years and for some variables. We find no evidence that CES accuracy has systematically improved over time, despite substantial changes to YouGov’s matching, weighting, and panel recruitment practices. In substantive terms, CES accuracy appears broadly comparable to—and in some cases worse than—what has been documented for other nonprobability survey designs.

% A central finding is that accuracy gains in the CES are increasingly driven by YouGov’s bespoke post-stratification weights rather than by sample matching alone. When CES data are analyzed without weights, error increases meaningfully, indicating that matching by itself does not reliably approximate a probability sample. Conversely, applying YouGov’s proprietary weights reduces error, particularly for variables directly incorporated into the post-stratification scheme. In contrast, best-practice calibration weighting using anesrake—while effective at aligning the sample to national benchmarks on the variables it constrains—does not systematically improve accuracy for outcomes outside the weighting set. This result replicates a broader pattern observed in prior work: traditional demographic post-stratification does not generalize well to unconstrained variables in nonprobability samples.

% These findings align with, and in some respects extend, a growing literature documenting the limits of nonprobability survey accuracy. Prior studies have shown that average error in nonprobability samples is often only modestly worse than in probability samples, but that nonprobability designs tend to exhibit substantially greater variance, including a nontrivial share of large errors. Our results are consistent with this pattern. Even when average RMSE appears moderate, the distribution of errors is highly skewed, with sizable deviations in particular years, offices, and states. In this sense, CES performance looks no better—and at times worse—than what has been reported for other nonprobability samples using simpler recruitment or weighting strategies. While YouGov’s approach is often described as exceptional, our results suggest that it remains subject to similar limitations to those other nonprobability sampling methods face.

% Nonprobability sampling is now a permanent feature of the survey research landscape. YouGov, in particular, occupies a dominant position in academic and applied research, and its use of sample matching has been widely interpreted as a methodological breakthrough. Our evidence challenges that interpretation. Sample matching, as implemented in the CES, does not appear to deliver accuracy gains that distinguish it meaningfully from other nonprobability methods. Instead, accuracy improvements depend heavily on post-stratification choices that incorporate outcome-adjacent information—an approach that may be infeasible for many research applications that lack benchmarks. For those cases, our secondary-variable results are the more relevant guide to expected accuracy.

% At the same time, our findings help reconcile apparent contradictions in the literature. Some prior evaluations of the CES report lower error in selected years or offices. We reproduce lower error estimates in some cases as well. However, we show that these conclusions are sensitive to which variables are evaluated and whether those variables were used in weighting. When considering the most informative outcomes—secondary variables that were not weighted and that appear consistently across years—accuracy remains persistently low and no evidence of improvement over time. This distinction between primary and secondary variables is crucial for interpreting CES performance and has been largely absent from prior evaluations.

% The first of several limitations of this study is that because YouGov’s weighting scheme has evolved over time, only a limited set of variables can be classified as consistently secondary across all survey waves, many of which involve candidate choice. Measurement of these outcomes is itself imperfect: special elections, fusion voting, and changes in ballot structure sometimes make it difficult to determine precisely which contest a respondent had in mind. These sources of noise likely inflate error estimates to some degree. That said, such limitations reflect real-world constraints faced by researchers using CES data and do not undermine the broader patterns we observe.

% Our results also point toward important directions for future research. First, the high variance of nonprobability samples implies that the single-year evaluations still common in the literature can be misleading. Multi-year designs are essential for assessing stability, trends, and tail risk in survey accuracy. Second, the sharp contrast between the effectiveness of weighting at the national level and its breakdown at the state level warrants closer investigation. Understanding why calibration succeeds for national margins but fails to generalize across subnational units may shed light on the underlying sources of error in nonprobability samples. Finally, comparative evaluations across multiple vendors remain rare. The Pew study remains one of the few head-to-head assessments of YouGov relative to other providers; additional work of this kind would be valuable.

% In sum, our findings challenge the narrative of sample matching as a final solution to the accuracy problems of nonprobability surveys. While matched samples combined with bespoke post-stratification can improve performance on selected variables, they do not escape the fundamental constraints that characterize nonprobability designs: substantial variance, sensitivity to weighting choices, and limited generalizability to unconstrained outcomes. Researchers relying on CES data should therefore interpret results with caution, particularly when drawing inferences about variables not explicitly targeted by the weighting process.

% ==========================================================================
%  APPENDIX A (printed with the article, before the references, as in POQ)
% ==========================================================================
\clearpage
\appendix

% Appendix A is the AAPOR disclosure appendix that POQ requires; its tables
% are numbered A1, A2, .... The online Supplementary Material follows the
% references, in sections A--F with tables, figures, and equations numbered
% S.1, S.2, ....

% ==========================================================================
%  APPENDIX A: DISCLOSURE ELEMENTS
% ==========================================================================

\section*{Appendix A. AAPOR-Required Disclosure Elements}
\phantomsection
\addcontentsline{toc}{section}{Appendix A. AAPOR-Required Disclosure Elements}
\label{app:disclosure}
\setcounter{table}{0}
\setcounter{figure}{0}
\renewcommand{\thetable}{A\arabic{table}}
\renewcommand{\thefigure}{A\arabic{figure}}
\renewcommand{\theHtable}{A\arabic{table}}
\renewcommand{\theHfigure}{A\arabic{figure}}

This appendix documents the study's data sources in the form required by the AAPOR Code, complementing their descriptions in ``The Cooperative Election Study'' and ``Benchmarks'' in the main text.

This study analyzes three existing data sources and collected no new data. For each source, the elements required by the AAPOR Code are given below. Where an element is described in permanently archived, study-specific documentation, we link to it; where an element is not documented or does not apply, we say so and why. Year-specific details for the CES are collected in table~\ref{tab:disclosure_by_year}.

\subsection*{Data Source 1: Cooperative Election Study (CES), 2006--2022}

\begin{description}[style=unboxed, leftmargin=0pt, font=\normalfont\bfseries]
\item[Data source:] Survey data. Secondary analysis of the publicly released CES Common Content for the nine federal election years 2006 through 2022. Data and documentation for each year are archived on Harvard Dataverse:
2006, \url{https://doi.org/10.7910/DVN/Q8HC9N};
2008, \url{https://doi.org/10.7910/DVN/YUYIVB};
2010, \url{https://doi.org/10.7910/DVN/VKKRWA};
2012, \url{https://doi.org/10.7910/DVN/HQEVPK};
2014, \url{https://doi.org/10.7910/DVN/XFXJVY};
2016, \url{https://doi.org/10.7910/DVN/GDF6Z0};
2018, \url{https://doi.org/10.7910/DVN/ZSBZ7K};
2020, \url{https://doi.org/10.7910/DVN/E9N6PH};
2022, \url{https://doi.org/10.7910/DVN/PR4L8P}.

\item[Data collection strategy:] Online surveys, with a pre-election and a post-election wave in each year. In 2006, additional profile information was collected in a separate survey in August 2006.

\item[Research sponsor and conductor:] The CES is sponsored by a consortium of universities and research teams, with support from the National Science Foundation. Data were collected by Polimetrix (2006), YouGov/Polimetrix (2008--2012), and YouGov (2014--2022). No funding was provided for this study.

\item[Measurement tools/instruments:] The exact wording, response options, and instructions for every question analyzed are in each year's questionnaire and guide, archived with the data above. Table~\ref{tab:ces_variable_map} lists, for each benchmarked item and year, the CES variable used, verified against the variable labels in the released data files. Wording and response options for several items differ across waves; each year's responses were recoded to the benchmark categories as described in ``Measuring Error.'' Vote-choice items asked respondents to choose among parties in earlier years and among named candidates, shown with party labels, in later years; see ``Measuring Error'' in the main text.

\item[Population under study:] The CES documentation describes its samples as representative of US citizens ages 18 and older, the population from which its sampling frames are drawn. Estimates are evaluated for the 50 states and the District of Columbia and, for US House, congressional districts, in each election year. Respondents who reported being non-citizen immigrants were excluded in every year (302 in 2006, 339 in 2008, 543 in 2010, 692 in 2012, 1,003 in 2014, 1,368 in 2016, 1,134 in 2018, 1,010 in 2020, and 624 in 2022), so CES estimates, like the CPS benchmarks (Data Source 2), refer to citizens; respondents who skipped the item were retained.

\item[Methods used to generate and recruit the sample:]
The samples are non-probability samples drawn from opt-in online panels: the Polimetrix/YouGov panel (PollingPoint through 2012) and, from 2008, panels supplied by other companies (tables~\ref{tab:externalproviders} and \ref{tab:disclosure_by_year}). Coverage is limited to people who joined those panels.

From 2008 on, each guide says a target sample was drawn from a sampling frame and the completed pre-election interviews were then matched to it (table~\ref{tab:ces_stages}). The 2006 wave followed a different procedure (Supplementary Material section E, ``The 2006 Wave''). From 2008 through 2016, frames were built from American Community Survey (ACS) citizen records, with political and religious variables attached from the Current Population Survey (CPS), exit polls, and the Pew Religious Landscape Survey. The combination of sources differed by year. The 2018 and 2020 frames were based on the 2017 and 2019 ACS, respectively, with voter registration matched from the CPS (the 2018 CPS in 2018; the 2018 and 2020 CPS in 2020). The guides for those years describe no other political or religious variables on the frame. In 2022, the frame was a ``politically representative modeled frame'' based on the 2019 ACS, with political variables modeled from voter files, the 2020 CPS Voting and Registration Supplement, the 2020 National Election Pool exit poll, and the 2020 CES.

Panelists were selected for invitation using a five-way cross-classification (age $\times$ gender $\times$ race $\times$ education $\times$ state) from 2008 through 2016, and a six-way cross-classification (age $\times$ gender $\times$ race $\times$ education $\times$ region $\times$ sample source) from 2018. No other quotas are documented, and the guides do not describe how the pool of panelists eligible for invitation was determined. \textcite{rivers_sample_2006} describes the general procedure YouGov developed for this purpose, in which invitations are allocated using a cross-classification on age, gender, race, education, region, and sample source, with the number of invitations per target case set from estimated completion probabilities so as to yield at least one completed interview per case. All pre-election respondents were re-invited to the post-election wave. The guides do not describe the incentives offered to panelists, except that in 2006 ``inducements'' were offered late in October to increase participation among low-income and low-interest respondents; YouGov currently describes its panel incentive as points redeemable for small rewards \parencite{yougov_methodology}, and incentives offered by external providers are not documented. Data collection was not in person.

\item[Method(s) and mode(s) of data collection:] Self-administered web questionnaires, administered in English.

\item[Dates of data collection:] Field periods for each year are in table~\ref{tab:sample_description}. Guides through 2014 report them only at monthly or relative precision.

\item[Sample sizes and discussion of precision:] Released sample sizes, the number of completed pre-election interviews in the matching pool, and post-election completes are in tables~\ref{tab:sample_description} and \ref{tab:disclosure_by_year}. The number of respondents underlying each estimate is shown in figure~\ref{fig:app_sample_size_distribution}, respondents per state in table~\ref{tab:states_respondents_by_office}, and the number of estimates by variable and year in table~\ref{tab:observation_counts}. Because the samples are not probability samples, we report no margins of sampling error; this study instead measures total error directly, as the RMSE of CES estimates against external benchmarks. The trend analysis (table~\ref{tab:rmse_trend_primary_secondary}) reports model-based HC3 standard errors and permutation $p$-values; these describe uncertainty in the fitted trend across variable-years, not sampling error in the CES estimates.

\item[Whether and how the data were weighted:] Estimates are analyzed unweighted (matching only), with the CES-provided weights, and with anesrake calibration weights we constructed. The CES weights were computed by weighting the matched sample to the frame (propensity-score subclassification in 2008--2012; entropy balancing from 2014, Supplementary Material section B) and then post-stratifying (raking in 2022), using the variables in table~\ref{tab:ces_stages}. The guides state that weights were trimmed at 7 (2008--2012) or 15 (2014--2022), but the released 2012 weight reaches 20 (Supplementary Material section F, item 1). Weights were normalized to the sample size (see ``CES Sample Matching and Weighting Procedure''). The CES weight variable used in each year is listed in table~\ref{tab:disclosure_by_year}. Where the CES provides a separate post-election weight (2016--2022), that weight is applied to every item, including demographics measured in the pre-election wave, so that all estimates in a year rest on the same respondents and weights. In 2014 the released file contains a single weight. The anesrake weights and their targets (CPS demographic benchmarks) are described in ``Weighting Scenarios'' and table~\ref{tab:vars_used_by_anesrake}. We used the R package \texttt{anesrake} with \texttt{type = "pctlim"}, \texttt{pctlim = 5}, and \texttt{choosemethod = "total"} and otherwise default settings. A variable was selected for raking if the absolute differences between its weighted CES category proportions and the CPS benchmark proportions, summed across its categories, totaled at least 5 percentage points; this rule decides which variables are used, not when raking stops. Selected variables were raked by iterative proportional fitting, in a fixed order, until the weights converged; any variable whose discrepancy then reached 5 points was added and raking was repeated. Weights were capped at 5 and re-centered to a mean of 1. Variables provided to \texttt{anesrake} but not selected are marked in table~\ref{tab:vars_used_by_anesrake}.

\item[How the data were processed and procedures to ensure data quality:] The guides state that pre-election interviews entered the matching pool ``after quality controls were applied'' but do not describe those controls. Our processing, including recoding CES responses to benchmark categories, is described in ``Measuring Error.'' Responses of ``don't know,'' ``not sure,'' ``did not vote in this race,'' and ``I did not vote,'' along with skipped, not-asked, and no-race responses, were set to missing and dropped item by item; ``other'' responses were retained in the denominator. No filter on reported turnout was applied to vote-choice items.

State-level estimates from 2008 onward use the state recorded in the post-election wave, so respondents who completed only the pre-election wave are excluded from all state-level estimates; the 2006 estimates use the pre-election state variable (\texttt{v1002}). US House estimates exclude respondents without a post-election congressional district; in 2020, a further 774 North Carolina respondents flagged in the CES data were excluded from House estimates. In 2018, respondents in eight states (Texas, Pennsylvania, Indiana, South Carolina, Alabama, Kentucky, Oklahoma, and Utah) who reported voting a straight party ticket (\texttt{CC18\_408} $= 2$; $n = 4{,}021$) were not asked about individual statewide and US House contests. Following the procedure described in the 2018 guide, we assigned each of these respondents the candidate of their ticket's party for US Senate, governor, US House, attorney general, and secretary of state; respondents whose ticket party fielded no candidate, who reported a minor-party ticket, or whose ticket party was not recorded were left missing. Their state legislative vote choices were taken directly from \texttt{CC18\_413c} and \texttt{CC18\_413d}. Also in 2018, the post-election questionnaire displayed Kansas's attorney general and secretary of state candidates to Kentucky respondents, where neither office was on the ballot, and displayed no candidates to Kansas respondents, none of whom has a recorded response to either item. The 2018 guide does not mention this error. Kentucky has no benchmark for these offices and Kansas has no responses, so both states are absent from the 2018 estimates for these offices.

CES candidate names were matched to official returns first by exact, case-insensitive comparison and otherwise by Jaro string distance (R package \texttt{stringdist} 0.9.15, method \texttt{"jw"} with $p = 0$), taking the closest candidate with no threshold. For offices other than president, when the benchmark candidate was a Democrat or Republican, distance matching considered only CES candidates of that party, and every match was checked for agreement with the party in the official returns. Of 4,185 candidate contests, 2,874 matched exactly and 1,309 by string distance (maximum distance 0.52); in the remaining two (US Senate in New Jersey, 2012 and 2018), Robert Menendez was matched to ``Bob Menendez'' by manual override. All 1,309 distance-based matches were reviewed manually and none required correction. Party affiliations were assigned manually where official returns omitted them (Supplementary Material section C).

\item[Panel description:] Panel recruitment is described in ``Panels and Invitations'' in the main text. Procedures for managing panel membership, participation, and attrition are not described in the CES guides.

\item[Interviewer details:] No interviewers were used; the questionnaires were self-administered. Human coding in this study was limited to the manual review of distance-based candidate matches, one candidate-name override, and party assignments described above, all performed by the authors.

\item[Screening criteria and process:] The frames were restricted to US citizens ages 18 and older. The guides do not document any other eligibility screening or whether respondents were routed from other surveys.

\item[Study stimuli:] None beyond the questionnaires archived with the data.

\item[Dispositions or response or participation rates:] Case dispositions and the rates the guides label AAPOR Response Rate 1, which divides completed interviews by all sampled cases not known to be ineligible, are reported in table~\ref{tab:disclosure_by_year} for 2008--2022; the 2006 guide reports none. The population of interest for an opt-in panel is not well defined \parencite[78]{aapor2023}, and the number of people exposed to a recruitment invitation is unknown. These figures are therefore completion rates among invited panelists, not response or participation rates for the population. Unlike a response rate in a probability sample, a high completion rate does not imply reduced risk of nonresponse bias. The 2014 and 2016 guides report YouGov and external-panel cases separately and combined, and the completion rates of external vendors were far lower (\GuideCompletionYouGov{} versus \GuideCompletionExternal{} in 2016). The 2018 guide says explicitly that its figures cover YouGov panelists only. The 2020 and 2022 guides do not, but they appear to, because their disposition tables count fewer completes (60,466 and 64,957) than the pooled YouGov and vendor interviews that entered matching (74,099 and 80,233). Even among these YouGov-only years, completion rates ranged from \CompletionRateRecentMin{} to \CompletionRateRecentMax{} percent, so the change in reporting base does not account for all of the variation. The 2008--2012 guides do not state which respondents the rate covers. Their sampling sections name external panels alongside YouGov's, and none of their disposition tables separates respondents by sample source. The 2008 table counts slightly fewer completes (49,043) than the roughly 50,800 pre-election interviews that entered matching, the same pattern as in 2020 and 2022; the 2010 guide gives no count of completed interviews entering matching, only a pool of ``over 150,000'' respondents; and the 2012 table is divided by study component rather than by sample source. A participation rate for the population cannot be computed.

\item[Measurement and model specification:] RMSE is defined in ``Measuring Error.'' The trend models are specified in table~\ref{tab:rmse_trend_primary_secondary} and figure~\ref{fig:rmse_trend_primary_secondary}, the sample-size model in figure~\ref{fig:respondent_count_sqrt_fit}, and the anesrake weights in ``Weighting Scenarios'' and under ``Whether and how the data were weighted'' above.

\item[Limitations:] The CES samples are non-probability samples, so their error cannot be characterized by sampling theory. Turnout, voting method, and vote choice are self-reports, so measured error includes misreporting as well as sample composition. Documentation is incomplete for some waves: the 2006 guide does not describe a weighting procedure, the 2012 guide repeats the 2010 description, several guides repeat identical figures for the size of the matching pool, and no guide describes the quality controls applied. Some benchmark outcomes (US Senate and gubernatorial results from 2014 onward, and presidential vote in 2020) were used in the CES weights, so weighted estimates for them are not independent tests of accuracy.
\end{description}

\begin{table}[htbp]
\centering
\caption{CES disclosure details by year.}
\label{tab:disclosure_by_year}
\footnotesize
\setlength{\tabcolsep}{4pt}
\begin{tabular}{@{}>{\raggedright\arraybackslash}p{0.7cm}>{\raggedright\arraybackslash}p{3.0cm}>{\raggedright\arraybackslash}p{4.7cm}>{\raggedright\arraybackslash}p{1.7cm}>{\raggedright\arraybackslash}p{1.6cm}>{\raggedright\arraybackslash}p{2.9cm}@{}}
\toprule
Year & Panels & Dispositions of invited cases (I / P / R / UH / NE) & RR1 & Post-election completes & CES weight used \\
\midrule
2006 & PollingPoint & Not reported & -- & 28,757 & \texttt{v1001} \\
2008 & PollingPoint; E-Rewards; Western Wats & 49,043 / 13,137 / 627 / 42,007 / 1,082 & .468 & 27,021 & \texttt{V201} \\
2010 & PollingPoint; E-Rewards; Western Wats & 75,450 / 27,155 / 4,645 / 79,723 / 9,262 & .404 & 46,684 & \texttt{V101} \\
2012 & PollingPoint; E-Rewards; Western Wats & 83,583 / 20,229 / 7,754 / 138,624 / 14,267$^a$ & .334$^a$ & 45,018 & \texttt{V103} \\
2014 & YouGov; MyPoints; Research Now; SSI & YouGov: 43,203 / 6,480 / 1,897 / 29,448 / 1,442 \newline External: 52,383 / 28,773 / 7,392 / 238,639 / 19,643 \newline (87,389 in matching pool) & .533 / .160 / .234$^b$ & 48,853 & \texttt{weight} \\
2016 & YouGov; MyPoints; Research Now; SSI; GMI & YouGov: 53,939 / 7,177 / 2,519 / 64,997 / 982 \newline External: 52,443 / 22,692 / 5,210 / 557,367 / 43,822 & .419 / .082 / .139$^b$ & 52,899 & \texttt{commonweight\_post} \\
2018 & YouGov; Dynata; Critical Mix; Prodege & YouGov only: 63,301 / 8,341 / 125 / 139,374 / -- & .300$^c$ & 51,808 & \texttt{commonpostweight} \\
2020 & YouGov; Dynata; Critical Mix; Prodege & 60,466 / 7,900 / 1,591 / 28,922 / 1,606 \newline (74,099 in matching pool) & .612 & 51,551 & \texttt{commonpostweight} \\
2022 & YouGov; Dynata; Disqo; Prodege; Generation Lab & 64,957 / 10,771 / 3,102 / 57,676 / 3,017 \newline (80,233 in matching pool) & .476 & 50,981 & \texttt{commonpostweight} \\
\bottomrule
\end{tabular}
\par\smallskip
\begin{minipage}{\linewidth}\footnotesize
\textit{Note:} I = complete; P = partial; R = refusal and breakoff; UH = unknown eligibility; NE = not eligible. Dispositions and RR1 are as reported in each year's CES guide; RR1 is reported in the main text as a completion rate. $^a$Common Content sample only; the 2012 guide also reports separate panel and NES samples. $^b$YouGov / external / combined. $^c$YouGov panelists only. The 2014, 2016, and 2018 guides each give the matching pool as ``approximately 87,389.'' We report the figure only for 2014, where it first appears, and treat the 2016 and 2018 figures as incorrectly carried forward without updating (Supplementary Material section F, item 3). Post-election completes are counts of respondents coded as completing the post-election wave (\texttt{v4001}, \texttt{V400}, \texttt{V403}, or \texttt{tookpost}). The 2012 file has no such variable, so its count is respondents with a non-missing post-election state (\texttt{inputstate\_post}).
\end{minipage}
\end{table}

\begin{table}[htbp]
\centering
\caption{CES variables used for each benchmarked item, by year.}
\label{tab:ces_variable_map}
\resizebox{\textwidth}{!}{%
\begin{tabular}{llllllllll}
\toprule
Item & 2006 & 2008 & 2010 & 2012 & 2014 & 2016 & 2018 & 2020 & 2022 \\
\midrule
Age Group & \texttt{v2020}$^a$ & \texttt{V207}$^a$ & \texttt{V207}$^a$ & \texttt{birthyr}$^a$ & \texttt{birthyr}$^a$ & \texttt{birthyr}$^a$ & \texttt{birthyr}$^a$ & \texttt{birthyr}$^a$ & \texttt{birthyr}$^a$ \\
Sex & \texttt{v2004} & \texttt{V208} & \texttt{V208} & \texttt{gender} & \texttt{gender} & \texttt{gender} & \texttt{gender} & \texttt{gender} & \texttt{gender4} \\
Education & \texttt{v2018} & \texttt{V213} & \texttt{V213} & \texttt{educ} & \texttt{educ} & \texttt{educ} & \texttt{educ} & \texttt{educ} & \texttt{educ} \\
Hispanic Origin & -- & -- & \texttt{V290} & \texttt{hispanic} & \texttt{hispanic} & \texttt{hispanic} & \texttt{hispanic} & \texttt{hispanic} & \texttt{hispanic} \\
Employment Status & \texttt{v2030} & \texttt{V209} & \texttt{V209} & \texttt{employ} & \texttt{employ} & \texttt{employ} & \texttt{employ} & \texttt{employ} & \texttt{employ} \\
Family Income & \texttt{v2032} & \texttt{V246} & \texttt{V246} & \texttt{faminc} & \texttt{faminc} & \texttt{faminc} & \texttt{faminc\_new} & \texttt{faminc\_new} & \texttt{faminc\_new} \\
Union Membership & \texttt{v2082} & \texttt{CC329} & \texttt{V264} & \texttt{union} & \texttt{union} & \texttt{union} & \texttt{union} & \texttt{union} & \texttt{union} \\
Veteran Status & \texttt{v3083} & \texttt{CC328\_3} & \texttt{V254} & \texttt{milstat\_3} & \texttt{milstat\_3} & \texttt{milstat\_3} & \texttt{milstat\_3} & \texttt{milstat\_3} & \texttt{milstat\_3} \\
Residence Duration & -- & \texttt{CC334} & \texttt{CC351} & \texttt{CC351} & \texttt{CC351} & \texttt{CC16\_361} & -- & \texttt{CC20\_361} & \texttt{CC22\_361} \\
Voting Turnout & \texttt{v4004} & \texttt{CC403} & \texttt{CC401} & \texttt{CC401} & \texttt{CC401} & \texttt{CC16\_401} & \texttt{CC18\_401} & \texttt{CC20\_401} & \texttt{CC22\_401} \\
Voting Method & \texttt{v4006} & \texttt{CC405} & \texttt{CC403} & \texttt{CC403} & \texttt{CC403} & \texttt{CC16\_403} & \texttt{CC18\_403} & \texttt{CC20\_403} & \texttt{CC22\_403} \\
President & -- & \texttt{CC410} & -- & \texttt{CC410a} & -- & \texttt{CC16\_410a} & -- & \texttt{CC20\_410} & -- \\
US Senate & \texttt{v4014} & \texttt{CC411} & \texttt{CC410a} & \texttt{CC410b} & \texttt{CC410b} & \texttt{CC16\_410b} & \texttt{CC18\_410b} & \texttt{CC20\_411} & \texttt{CC22\_411} \\
US House & \texttt{v4015} & \texttt{CC412} & \texttt{CC412} & \texttt{CC412} & \texttt{CC412} & \texttt{CC16\_412} & \texttt{CC18\_412} & \texttt{CC20\_412} & \texttt{CC22\_412} \\
Governor & \texttt{v4013} & \texttt{CC413} & \texttt{CC411} & \texttt{CC411} & \texttt{CC411} & \texttt{CC16\_411} & \texttt{CC18\_411} & \texttt{CC20\_413} & \texttt{CC22\_413} \\
Attorney General & \texttt{v4017} & -- & \texttt{CC413a} & \texttt{CC413a} & \texttt{CC413a} & \texttt{CC16\_413a} & \texttt{CC18\_420a} & \texttt{CC20\_414a} & \texttt{CC22\_414a} \\
Secretary of State & \texttt{v4018} & -- & \texttt{CC413b} & \texttt{CC413b} & \texttt{CC413b} & \texttt{CC16\_413b} & \texttt{CC18\_420b} & \texttt{CC20\_414b} & \texttt{CC22\_414b} \\
State Senator & \texttt{v4021} & \texttt{CC414\_2} & \texttt{CC413c} & \texttt{CC413c} & \texttt{CC413c} & \texttt{CC16\_413c} & \texttt{CC18\_413c} & \texttt{CC20\_415c} & \texttt{CC22\_415c} \\
State Representative & \texttt{v4020} & \texttt{CC414\_1} & \texttt{CC413d} & \texttt{CC413d} & \texttt{CC413d} & \texttt{CC16\_413d} & \texttt{CC18\_413d} & \texttt{CC20\_415d} & \texttt{CC22\_415d} \\
\midrule
Citizenship & \texttt{v3086} & \texttt{CC332} & \texttt{V263} & \texttt{immstat} & \texttt{immstat} & \texttt{immstat} & \texttt{immstat} & \texttt{immstat} & \texttt{immstat} \\
State assignment & \texttt{v1002}$^b$ & \texttt{V259} & \texttt{V206\_post} & \texttt{inputstate\_post} & \texttt{inputstate\_post} & \texttt{inputstate\_post} & \texttt{inputstate\_post} & \texttt{inputstate\_post} & \texttt{inputstate\_post} \\
Congressional district & \texttt{v1003} & \texttt{V264} & \texttt{V276\_post} & \texttt{cdid113\_post} & \texttt{cdid\_post} & \texttt{cdid115\_post} & \texttt{cdid116\_post} & \texttt{cdid116\_post} & \texttt{cdid118\_post} \\
\bottomrule
\end{tabular}}
\par\smallskip
\begin{minipage}{\linewidth}\footnotesize
\textit{Note:} Variable names are as they appear in each year's released Common Content file and were verified against the variable labels in those files; question wording is in the corresponding guide and questionnaire on Dataverse. Dashes indicate that the item was not measured, or the office was not on the ballot, that year. Vote-choice items are from the post-election wave. $^a$Age group is derived from year of birth. $^b$The 2006 file has no post-election state variable; the pre-election state is used.
\end{minipage}
\end{table}

\subsection*{Data Source 2: Current Population Survey (CPS)}

\begin{description}[style=unboxed, leftmargin=0pt, font=\normalfont\bfseries]
\item[Data source:] Survey data. Secondary analysis of the CPS, obtained through IPUMS CPS \parencite{flood_ipums_2024}, used for demographic and voting-method benchmarks.

\item[Data collection strategy, sponsor, and conductor:] A monthly household survey sponsored by the US Census Bureau and the Bureau of Labor Statistics and conducted by the Census Bureau. Methodology, questionnaires, and sampling are documented at \url{https://www.bls.gov/cps/documentation.htm} and \url{https://cps.ipums.org/cps/}.

\item[Measurement tools/instruments:] Benchmarks come from the CPS basic monthly files, with all twelve months pooled within each year (the March file used is the basic monthly March sample, not the Annual Social and Economic Supplement), obtained through the IPUMS API (IPUMS CPS version 13.0). Sex, education, age, Hispanic origin, veteran status, employment status, and family income are from the basic monthly items; union membership is from the earnings (outgoing rotation) items; residence duration and voting method are from the November Voting and Registration Supplement. Questionnaires are documented at the links above.

\item[Population, sample, mode, and dates:] A monthly probability sample of US households conducted by the US Census Bureau since 1940. Sampled households have been interviewed once a month for four months, then given eight months off, and again once a month for four additional months. Initial interviews were conducted in person, and subsequent interviews were conducted in person or by telephone. Interviewers attempted to collect data from all household members and accepted proxy reports from one household member about another household member when necessary. Data from 2006 through 2022 were used. Benchmarks are restricted to respondents ages 18 and older who are US citizens (IPUMS \texttt{CITIZEN} codes 1--4).

\item[Sample sizes, response rates, and weighting:] Annual sample sizes and household response rates are in table~\ref{tab:cps_sample_description}; sample sizes count person-month records, not unique persons. Benchmarks use the CPS weights (\texttt{WTFINL} for basic monthly items, \texttt{EARNWT} for union membership, and \texttt{VOSUPPWT} for Voting and Registration Supplement items), which account for sampling design, nonresponse, and undercoverage using population controls derived from the decennial census and updated annually.

\item[Other elements:] Interviewer procedures, screening, processing, and panel rotation are documented by the Census Bureau and BLS at the links above.

\item[Limitations:] CPS benchmarks carry their own sampling and nonsampling error, including proxy reports and, for voting method, self-report error; they are treated here as fixed, so part of the measured CES error may reflect benchmark error, particularly in small states. Race was not evaluated because CPS race categories could not be mapped exhaustively onto CES categories; Hispanic origin, measured comparably in both, was evaluated from 2010 onward.
\end{description}

\subsection*{Data Source 3: Election Returns and Turnout}

\begin{description}[style=unboxed, leftmargin=0pt, font=\normalfont\bfseries]
\item[Data source:] Government records. Official election returns, compiled by the MIT Election Data and Science Lab \parencite{mit_election_data_and_science_lab_us_2017-1}, the Federal Election Commission, CQ Press, Ballotpedia, and state election offices, and turnout from Michael McDonald's turnout rates, distributed by the University of Florida Election Lab \parencite{mcdonald_voter_nodate} (table~\ref{tab:data_sources}).

\item[Measurement:] Turnout is total ballots counted divided by the voting-eligible population (VEP) with ineligible felons added back, from McDonald's 1980--2022 General Election Turnout Rates, version 1.2 \parencite{mcdonald_voter_nodate}. This denominator is McDonald's estimate of the citizen voting-age population: it excludes noncitizens but, unlike the standard VEP, retains people disenfranchised by a felony conviction, matching the CES population of adult citizens. In 28 state-years for which total ballots are not reported, the numerator is votes cast for the highest office on the ballot: Alabama 2006 and 2012; Connecticut 2008; Delaware 2012; Maine 2006; Mississippi 2006--2016; Missouri 2012; New Mexico 2014; Oklahoma 2006, 2012, and 2016; Pennsylvania 2006, 2012, and 2016; Rhode Island 2012; Texas 2006, 2008, 2012, and 2014; Virginia 2016; West Virginia 2012 and 2016; and Wisconsin 2016. National turnout is summed ballots divided by the summed denominator across the 50 states and the District of Columbia. State-level turnout benchmarks cover the 50 states; the District of Columbia is excluded. Vote shares are each candidate's share of all votes cast in the contest. Each contest is compared on the candidate who received the most votes, except in eight contests, listed in Supplementary Material section C, where the second-place major-party nominee is used instead.

\item[Processing:] Matching of candidates and parties to CES responses, and the handling of multi-party endorsements and special elections, are described in Supplementary Material section C.

\item[Elements that do not apply:] These are official counts rather than samples, so sampling, mode, response rates, weighting, interviewers, and sampling precision do not apply.
\end{description}

\input{output/data_sources}                           % Table A3: tab:data_sources

\clearpage

% ==========================================================================
%  END MATTER (POQ order: Supplementary Material, Funding, Data Availability,
%  then References)
% ==========================================================================
\section*{Supplementary Material}
Supplementary Material sections A--F follow the references.

\section*{Funding}
No funding was provided for this study.

\section*{Data Availability}
The CES data analyzed in this study are publicly available on Harvard Dataverse; Appendix A gives the archive for each year. Replication data and code will be deposited in POQ's Dataverse upon acceptance.

\begin{singlespace}
\setlength{\bibitemsep}{0.5\baselineskip}
\printbibliography
\end{singlespace}

\clearpage
% ==========================================================================
%  SUPPLEMENTARY MATERIAL (online), sections A--F
% ==========================================================================
\section*{Online Supplementary Material}
\phantomsection
\addcontentsline{toc}{section}{Online Supplementary Material}
\setcounter{table}{0}
\setcounter{figure}{0}
\setcounter{equation}{0}
\renewcommand{\thetable}{S.\arabic{table}}
\renewcommand{\thefigure}{S.\arabic{figure}}
\renewcommand{\theequation}{S.\arabic{equation}}
\renewcommand{\theHtable}{S.\arabic{table}}
\renewcommand{\theHfigure}{S.\arabic{figure}}
\renewcommand{\theHequation}{S.\arabic{equation}}

% ==========================================================================
%  SECTION A: SUPPLEMENTARY TABLES AND FIGURES
% ==========================================================================
\subsection*{Section A. Supplementary Tables and Figures}
\phantomsection
\addcontentsline{toc}{subsection}{Section A. Supplementary Tables and Figures}

This section collects supplementary tables and figures that support the main text. Table~\ref{tab:national_anesrake_variables_accuracy} supports ``Effect of Conventional Post-Stratification on CES Error,'' reporting how closely the anesrake weights reproduced the CPS benchmarks used to build them. Table~\ref{tab:cps_sample_description} reports the CPS sample sizes and response rates behind the demographic benchmarks described in ``Benchmarks'' and in Appendix A's CPS disclosure. Table~\ref{tab:demographic_candidate_choice_weighting_comparison_vertical} and figure~\ref{fig:demographic_candidate_choice_weighting_comparison_vertical} compare how weighting affected primary and secondary demographic variables and the two candidate-choice offices, Governor and US Senate, that switched status in 2014. The remaining two figures extend the sample-size analysis in ``Sample Size and Estimation Error'': figure~\ref{fig:app_sample_size_distribution} reports the distribution of respondent counts underlying every error estimate; about two-thirds rest on 500 or fewer respondents, though the range is wide. Figure~\ref{fig:sample_size_distribution_rmse_state_year_pools} reports the corresponding distribution, and RMSE, at the level of state-year respondent pools rather than individual estimates.

% ---- Supplementary tables ------------------------------------------------
% Cite each table by the \label inside its own file, e.g.
% Table~\ref{tab:us_house_rmse_grid}. Do not add bare \label commands here:
% after an unnumbered heading they pick up the last section number, which is
% how the body text came to read "Appendix 4".
\input{output/national_anesrake_variables_accuracy}   % Table S.1: tab:national_anesrake_variables_accuracy
\input{output/cps_sample_description}                 % Table S.2: tab:cps_sample_description
\input{output/rmse_demographic_candidate_choice_weighting_comparison_vertical} % Table S.3

% ---- Figure S.1 -------------------------------------------------------------
\begin{figure}[!htbp]
    \centering
    \includegraphics[width=\linewidth, height=0.66\textheight, keepaspectratio]{output/rmse_demographic_candidate_choice_weighting_comparison_vertical.png}
    \caption[Weighting and RMSE by variable type]{Weighting and RMSE by variable type. Bars compare matching-only and weighted RMSE. $\Delta$ is weighted minus matching-only RMSE, so negative values indicate improvement. Panels A and B contain primary demographics (Age Group, Education, and Sex) and secondary demographics (Employment Status, Family Income, Union Membership, and Veteran Status), respectively. Panels C and D contain Governor and US Senate, which were secondary in 2008--2012 and primary in 2014--2022; blank years indicate the other classification. Weighting improves primary demographics more than secondary demographics and produces larger reductions for the candidate-choice variables during their primary years, although the latter comparison is confounded with time.}
    \par\smallskip
    \noindent\begin{minipage}{\linewidth}\footnotesize\textit{Alt text:} Four panels of paired bars comparing error before and after weighting, 2008 to 2022. Weighting lowered error in every year for primary demographics, Governor, and United States Senate, but changed it little for secondary demographics.\end{minipage}
    \label{fig:demographic_candidate_choice_weighting_comparison_vertical}
\end{figure}

% ---- Supplementary sample-size figures ----------------------------------

\begin{figure}[!htbp]
    \centering
    \includegraphics[width=\linewidth]{output/ces_error_sample_size_distribution_variable_width_frequency.png}
    \caption[Distribution of sample sizes underlying CES error estimates]{Distribution of sample sizes underlying CES error estimates. Sample size is the number of CES respondents contributing to the survey estimate whose error is measured against an external benchmark. For statewide estimates, this is the relevant state--year respondent count; for US House, which is a district-specific estimate, it is the relevant district--year respondent count. Bars represent unequal sample-size intervals: their heights report the number of estimate-level observations in each interval, their widths reflect the range of sample sizes covered, and their labels report the percentage of all included estimates. The figure includes all \NObsTotal{} unweighted CES estimates.}
    \par\smallskip
    \noindent\begin{minipage}{\linewidth}\footnotesize\textit{Alt text:} Histogram of error estimates by sample size, with an inset for ranges up to 500 respondents. About two-thirds rest on 500 or fewer respondents, and the 51 to 100 range holds \FigBTwoShareFiftyToHundred{} percent.\end{minipage}
    \label{fig:app_sample_size_distribution}   % Figure S.2
\end{figure}

\begin{figure}[!htbp]
    \centering
    \includegraphics[width=\linewidth]{output/sample_size_distribution_rmse_state_year_pools.png}
    \caption[Sample-size distribution and RMSE across state-year respondent pools]{Sample-size distribution and RMSE across state-year respondent pools. This figure provides a sample-level counterpart to the estimate-level sample-size distribution reported in figure~\ref{fig:app_sample_size_distribution}. In that figure, each variable-specific comparison with an external benchmark is an error observation. Consequently, the same underlying state-year respondent pool appears repeatedly because its respondents are used to estimate multiple variables and, where applicable, multiple district-specific outcomes. Here, each state-year respondent pool appears only once. Its sample size is defined as the largest valid respondent count observed for that state-year; smaller variable-specific counts may result from item nonresponse or geographically restricted estimates. Bar heights report the percentage of state-year pools in each sample-size range, and labels report the number and percentage of pools. RMSE is calculated hierarchically so that variables producing more observations and pools producing more error rows do not receive disproportionate weight. Squared percentage-point errors are first averaged within each state-year-variable combination. Those variable-level mean squared errors are then averaged within each state-year pool, giving every variable equal weight. Within each sample-size range, the pool-level mean squared errors are averaged across state-year pools and the square root is taken, giving every pool equal weight. The overall equal-pool RMSE shown in the figure is calculated analogously across all state-year pools. Results use unweighted CES estimates and exclude turnout to match the corresponding main-text analysis.}
    \par\smallskip
    \noindent\begin{minipage}{\linewidth}\footnotesize\textit{Alt text:} Bar chart of \NStateYearPools{} state and year respondent pools by sample size, with a line for error. Error is about \PoolRmseSmallest{} percentage points for the smallest pools and \PoolRmseAboveHundredMin{} to \PoolRmseAboveHundredMax{} above 100 respondents, with no steady decline.\end{minipage}
    \label{fig:sample_size_distribution_rmse_state_year_pools}

    
\end{figure}

% Place all section A floats before section B begins.
\clearpage

% ---- Text sections ---------------------------------------------------------
% Unnumbered headings. \ref to these labels would print a stale section
% number, and pandoc drops \nameref, so refer to them by name in the text.
% The labels stay as link anchors (they become the Word bookmarks).

\subsection*{Section B. Entropy Balancing Procedure}
\phantomsection
\addcontentsline{toc}{subsection}{Section B. Entropy Balancing Procedure}
\label{app:entropy_balancing}

This section reconstructs the entropy-balancing procedure used in the first weighting step described in ``CES Sample Matching and Weighting Procedure,'' since the CES guides report the method's balancing conditions only in general terms.

Beginning in 2014, the Cooperative Election Study (CES) weighted the matched sample to the sampling frame using entropy balancing \parencite{hainmueller_entropy_2012}. Entropy balancing is a multivariate reweighting method that assigns scalar weights to individual observations such that the reweighted sample exactly satisfies a set of covariate balance constraints while remaining as close as possible to a set of base weights.

Formally, weights were obtained by minimizing Kullback--Leibler divergence \parencite{kullback_information_1959} between the final weights \( w_i \) and base weights \( q_i \):
\[
\min H(w) = \sum_{i} w_{i} \log\!\left(\frac{w_{i}}{q_{i}}\right),
\]
subject to the following constraints:
\begin{gather}
    \sum_{i \in S} w_{i} c_{ri}(X_i) = 0, \label{eq:ebalance1} \\
    \sum_{i \in S} w_{i} = 1, \label{eq:ebalance2} \\
    w_{i} \ge 0 \quad \forall i. \label{eq:ebalance3}
\end{gather}

Here, \( S \) indexes the matched sample, \( w_i \) denotes the entropy-balanced weight assigned to respondent \( i \), and \( q_i \) denotes the base weight. In the standard formulation, base weights are uniform, \( q_i = 1/N \), where \( N \) is the matched sample size \parencite{hainmueller_entropy_2012}.

Constraint~\eqref{eq:ebalance1} requires the weighted sample to reproduce the frame value of the $r$-th balance condition. The CES guides for 2014--2022 report that first-stage weights were obtained by entropy balancing the matched cases to the sampling frame on ``multiple moment conditions'' in a listed set of covariates ``plus their interactions,'' but they do not report the functional form of those conditions. Following \textcite{hainmueller_entropy_2012}, we assume first-moment constraints of the form
\[
c_{ri}(X_i) = X_{ri} - \mu_r ,
\]
where $X_{ri}$ is the value of the $r$-th balancing term for respondent $i$ (a main-effect indicator or an interaction of indicators) and $\mu_r$ is its mean in the sampling frame. With categorical covariates coded as indicators, these constraints reproduce the frame's marginal and joint proportions, consistent with the 2018--2022 guides' description of weighting ``to match the distributions'' of the frame. The guides do not indicate whether higher-order moments were also balanced. Constraints~\eqref{eq:ebalance2} and~\eqref{eq:ebalance3} ensure that the resulting weights form a valid probability distribution.

Minimizing Kullback--Leibler divergence subject to these constraints yields a set of weights that deviate minimally from the base weights while achieving exact balance on the specified covariates. This property ensures that weights vary smoothly across observations and preserves their interpretation as probability weights.


\subsection*{Section C. Inconsistencies and Intricacies of Election Information Reporting}
\phantomsection
\addcontentsline{toc}{subsection}{Section C. Inconsistencies and Intricacies of Election Information Reporting}
\label{app:inconsistencies_election_info}

This section details the judgment calls involved in matching CES-reported vote choice to official election returns, which underlie the candidate-level comparisons described in ``Measuring Error'' and are cited in Appendix A's disclosure for the election-returns data. It covers candidate-name matching, the assignment of unaffiliated candidates to a party, the aggregation of multiply-endorsed candidacies, the treatment of simultaneous special and general elections, and the eight contests benchmarked against a major-party nominee rather than the top vote-getter.

The change in the wording of the CES vote-choice question, and the comparison rule that followed from it, are described in ``Measuring Error.'' Party-level reporting could not be cleanly mapped to survey responses in contests with multiple candidates sharing a party label, candidates appearing on multiple party lines, ballots without party labels, and candidates switching party affiliation between the primary and general election. Candidate vote choice was available for all presidential contests in all years, for US Senate and House elections beginning in 2008, for all gubernatorial elections beginning in 2010, and for attorney general and secretary of state races beginning in 2018. State legislative contests were evaluated using party vote share in all years.

Candidate names occasionally differed in format between the CES questionnaire and official election returns. Jaro string distance \parencite{jaro_advances_1989} was used to map CES candidates to those in the election returns; where no exact match existed, the closest name was taken, and all such matches were manually reviewed.

Some states did not report candidate party affiliation in official election returns, even when candidates were formally affiliated with a party. In these cases, candidates were manually assigned to their affiliated party to avoid artificially inflating CES error, as such candidates would otherwise have been treated as unaffiliated despite being identified by CES respondents as Democrats or Republicans.

In states such as New York and Vermont, candidates endorsed by multiple parties appeared multiple times on the ballot—once for each party line. For example, a candidate may have received votes under the Democratic Party, Working Families Party, and Progressive Party. In these cases, all votes in the election returns were aggregated to the candidate’s major-party affiliation (Democratic or Republican).

If the only election for an office in a year was a special election, that special election was used as the contest. If a state held both a regular and a special election for the same office on the same day, as happened for US Senate, the special election was excluded and only the regular election was used, because the CES typically asks about a single race for that office and it is ambiguous which one respondents reported.

In eight contests, CES vote choice is compared with the benchmark share of a major-party nominee rather than that of the top vote-getter: US Senate in Alaska (2010), Maine (2012 and 2018), Vermont (2012 and 2018), and Louisiana (2016); governor of Rhode Island (2010); and secretary of state of North Dakota (2018). In all but Louisiana, the winner was an independent or write-in candidate; the major-party nominee is used so that these candidate-level comparisons align with the party-level comparisons, in which independent and write-in votes fall into the ``Other'' category. In Louisiana, the leading candidate was not offered in the CES item. In each of the eight contests, the nominee used finished second.


\subsection*{Section D. US House Sample Size and Competitiveness}
\phantomsection
\addcontentsline{toc}{subsection}{Section D. US House Sample Size and Competitiveness}
\label{app:us_house_grid}

Table~\ref{tab:us_house_rmse_grid} supplements the main-text section ``Error as a Function of Contest Competitiveness'' by using US House races. House estimates are made by congressional district, and districts vary in how many CES respondents they contain, so these races let us examine sample size and competitiveness together. Columns give ranges of the winner's share of the vote. Rows give a minimum number of respondents per district, and each row keeps only district-years at or above that threshold. The upper panel reports RMSE in percentage points for the district-years in each cell, and the lower panel reports the number of district-years in each cell, both pooled over 2006--2022 and estimated with matching plus the CES weights.

\input{output/us_house_rmse_grid}                     % Table S.4: tab:us_house_rmse_grid

% Keep table S.4 inside section D.
\clearpage


% ==========================================================================
%  SECTION E: YEAR-SPECIFIC DETAILS OF THE CES PROCEDURE
% ==========================================================================
\clearpage
\subsection*{Section E. Year-Specific Details of the CES Sampling and Weighting Procedure}
\phantomsection
\addcontentsline{toc}{subsection}{Section E. Year-Specific Details of the CES Sampling and Weighting Procedure}
\label{app:ces_procedure_by_year}

This section records the year-to-year differences in the procedure described in ``CES Sample Matching and Weighting Procedure,'' as documented in each year's CES guide. Invitation cross-classifications, trimming thresholds, and the weight variable used in each year are given in Appendix~A, and the variables used at each stage in table~\ref{tab:ces_stages}; they are not repeated here.

\subsubsection*{The 2006 Wave}
The 2006 study, conducted by Polimetrix before its acquisition by YouGov, followed a different order from later waves. For general-population studies, the 2006 guide describes the target population as adults enumerated in consumer databases maintained by commercial vendors (Acxiom, Experian, and InfoUSA). Later waves used the ACS. The target sample was stratified by registration status, state size, and district competitiveness (16 strata), and further by age, race, and gender. Panelists came from a pool of more than 150,000 PollingPoint panelists, and selection had two steps. First, before fielding, several panelists were matched to each target case, until the expected number of responses per case was at least one, using a model of each panelist's probability of responding. Second, when a panelist began the interview, they were assigned to their own target case if it was still open, and otherwise to the closest open target case across all studies then in the field. The matching was thus set up before data collection, but final assignment happened as panelists began the interview. Cases were not matched to a pool of completed interviews. The 2006 guide does not list the matching variables or describe how the released weight (\texttt{v1001}) was built.

\subsubsection*{Sampling Frame}
Table~\ref{tab:frame_sources_by_year} lists, for each year from 2008, the ACS file from which the frame was drawn and the sources from which political and religious variables were attached. For the CPS data in 2008--2016, the guides state that the data were attached ``using a weighted Euclidean distance metric''; for the other sources, they say only that the data were ``matched.'' The 2022 guide instead describes a ``modeled frame,'' whose political variables were modeled from voter files, the 2020 CPS Voting and Registration Supplement, the 2020 National Election Pool exit poll, and the 2020 CES.

Some of these frame descriptions appear to have been carried forward from earlier guides (section F, items 1 and 12).

\begin{table}[!htbp]
\centering
\footnotesize
\begin{threeparttable}
\caption{Sampling frame sources by year, as described in the CES guides.}
\label{tab:frame_sources_by_year}
\begin{tabular}{l >{\raggedright\arraybackslash}p{2.6cm} >{\raggedright\arraybackslash}p{4.4cm} >{\raggedright\arraybackslash}p{5.2cm}}
\toprule
Year & ACS file & Registration and turnout & Other attached variables \\
\midrule
2008 & 2006 & November 2004 CPS & 2007 Pew Religious Landscape Survey$^a$ \\
2010 & 2008 & November 2008 CPS & 2007 Pew Religious Landscape Survey$^a$ \\
2012 & 2008$^b$ & November 2008 CPS$^b$ & 2007 Pew Religious Landscape Survey$^a$ \\
2014 & 2010 & November 2010 and 2008 CPS & 2007 Pew Religious Landscape Survey$^a$ \\
2016 & 2012 and 2010$^c$ & 2012 CPS (2012 frame); November 2010 and 2008 CPS (2010 frame) & Vote choice from 2012 national and state exit polls (2012 frame); 2007 Pew Religious Landscape Survey$^a$ (2010 frame) \\
2018 & 2017 & 2018 CPS (registration only) & None described \\
2020 & 2019 & 2018 and 2020 CPS (registration only) & None described \\
2022 & 2019 & Modeled$^d$ & Modeled$^d$ \\
\bottomrule
\end{tabular}
\begin{tablenotes}\footnotesize
\item \textit{Note:} $^a$Religion, church attendance, born-again or evangelical status, news interest, party identification, and ideology. $^b$Identical to the 2010 guide. $^c$Two frames; the guide does not describe how they were combined. $^d$From voter files, the 2020 CPS Voting and Registration Supplement, the 2020 National Election Pool exit poll, and the 2020 CES.
\end{tablenotes}
\end{threeparttable}
\end{table}

\subsubsection*{Target Sample}
From 2008 through 2012, the target sample was stratified by age, race, gender, and education; in 2014 and 2016, voter registration was added. The 2018--2022 guides state that a target sample was drawn and that completed interviews were ``matched to the target frame,'' but do not describe how the target sample was drawn (section F, item 11).

\subsubsection*{Selection from the Pool}
The guides for 2008 through 2016 print the distance function in full. The 2008, 2010, and 2012 functions are identical and combine gender, age, race, region, education, political interest, marital status, party identification, ideology, religion, church attendance, family income, voter registration, and metropolitan residence. The 2014 and 2016 functions combine age, race, education, gender, employment status, ideology, party identification measured in a baseline survey, born-again status, and voter registration, with interaction terms (several in 2014, one in 2016) and different weights in the two years. Both also include whether the respondent completed the post-election questionnaire, which implies that selection in those years took place after the post-election wave had been fielded (section F, item 14). The guides' lists of matching variables for those years also include political interest (section F, item 4). The 2018--2022 guides do not print a function; they describe selection as ``conditioning on registration status $\times$ age $\times$ race $\times$ gender $\times$ education.'' The guides call the metric a ``weighted Euclidean distance,'' but the printed functions are weighted sums of absolute differences and mismatch indicators, the form given in the main text (section F, item 16).

\subsubsection*{Weighting to the Frame}
In 2014 and 2016, the entropy-balancing conditions were age, gender, education, race, voter registration, ideology, baseline party identification, born-again status, and political interest, with their interactions. In 2018--2022, the guides list age, gender, education, race, and Hispanic origin as the frame distributions the sample was weighted to, plus 2020 presidential vote and turnout in 2022. The moment conditions they list omit Hispanic origin and, in 2022, turnout (section F, item 5). The 2018--2022 guides are unclear about whether the completed cases or the matched cases were weighted to the frame, and we assume the matched cases, as in earlier years (section F, item 10).

\subsubsection*{Post-Stratification}
From 2008 through 2012, the frame-weighted sample was post-stratified by gender, race, education, and age; in 2014 and 2016, voter registration was added. In 2018, the variables were age, gender, education, race, born-again status, voter registration, and 2016 presidential vote, and in 2020 the same variables plus 2020 presidential vote. In 2022, the weights were instead raked (iterative proportional fitting) on age, gender, education, race, born-again status, voter registration, 2020 presidential vote, and some joint distributions of these. The guides' descriptions of the additional state-level adjustment, and the offices it covered, are discussed in section F, item 8.

\subsubsection*{Post-Election Weights}
The 2016 file includes pre-election and post-election weights, each in versions computed before and after vote validation. From 2018, the files include \texttt{commonweight} for all respondents and \texttt{commonpostweight} for respondents who completed both questionnaires, along with weights for respondents matched to a validated registration record. The guides describe the target population of each weight but not how the post-election weights were computed (section F, item 17). The weight used in each year of this study is listed in table~\ref{tab:disclosure_by_year}.

% ==========================================================================
%  SECTION F: PROBLEMS IN THE CES GUIDES' SAMPLING AND WEIGHTING DOCUMENTATION
% ==========================================================================
% Screenshots of the guides are in guide_evidence/ (tables_and_figures/guide_evidence
% in the repository); they are cited in the text as Evidence 1--12.
\newcommand{\guideevidence}[3]{%
  \par\medskip\noindent
  \begin{minipage}{\linewidth}\centering
  \includegraphics[width=#2]{guide_evidence/#1}\par\smallskip
  \begin{singlespace}\raggedright\small\textit{#3}\par\end{singlespace}
  \end{minipage}\par\medskip}

\clearpage
\subsection*{Section F. Problems in the CES Guides' Sampling and Weighting Documentation}
\phantomsection
\addcontentsline{toc}{subsection}{Section F. Problems in the CES Guides' Sampling and Weighting Documentation}
\label{app:guide_documentation_problems}

This section describes instances in which the CES guides' sampling-and-weighting documentation appears to be incorrect or incomplete, across the eight waves from 2008 to 2022. For each item, we offer the evidence of a possible problem. Table~\ref{tab:guide_problems_summary} summarizes the items.

Page numbers refer to the PDF pages of each CES guide.

\begin{table}[!htbp]
\centering
\footnotesize
\begin{threeparttable}
\caption{Possible problems in the CES guides' sampling and weighting documentation.}
\label{tab:guide_problems_summary}
\begin{tabular}{l >{\raggedright\arraybackslash}p{10.2cm} >{\raggedright\arraybackslash}p{2.4cm}}
\toprule
Item & Problem & Years affected \\
\midrule
1 & The 2012 sampling-and-weighting section is a copy of the 2010 sampling-and-weighting section with no revision. & 2012 \\
2 & The 2014 codebook entry for the common content weight is identical to the 2012 codebook entry for the common content weight. & 2014 \\
3 & The 2014 pool figure ``approximately 87,389'' is repeated in the 2016 and 2018 guides. & 2016, 2018 \\
4 & The variables listed as used for matching disagree with the variables that appear in the reported distance functions. & 2014, 2016 \\
5 & Hispanic origin (and, in 2022, turnout) is listed as a weighting target but not as a moment condition. & 2018--2022 \\
6 & The interaction terms in the moment conditions are not specified. & 2014--2022 \\
7 & Turnout is added to the 2014 frame but not used in any later step. & 2014 \\
8 & The offices in the statewide-races adjustment were not named. & 2014, 2016 \\
9 & No source is given for the born-again target. & 2018--2022 \\
10 & Which cases were weighted to the frame is unclear. & 2018--2022 \\
11 & The target-sample strata were not described. & 2018--2022 \\
12 & How the two frames were combined is never described, and the second frame is not described at all in 2018 and 2020. & 2016, 2018, 2020 \\
13 & How the modeled frame was built is not described. & 2022 \\
14 & The printed distance functions include completion of the post-election wave, which implies selection after the election, although the guides describe matching completed pre-election interviews. & 2014, 2016 \\
15 & The 2018 and 2020 guides cite CPS registration data that were released after each election. & 2018, 2020 \\
16 & The distance metric is mislabeled as a weighted Euclidean distance. & 2008--2016 \\
17 & How the post-election weights were computed is not described. & 2016--2022 \\
\bottomrule
\end{tabular}
\end{threeparttable}
\end{table}

% --------------------------------------------------------------------------
\subsubsection*{Text Possibly Carried Forward from an Earlier Guide without Being Updated}

For each issue below, duplicated text or statistics suggest that the description was copied from an earlier guide without being updated, which raises the possibility that key information, such as which variables were used to construct the weights, might be incorrect.

\paragraph*{1. The 2012 sampling-and-weighting section is a copy of the 2010 section (year with possible errors: 2012)}
The 2012 guide's sampling-and-weighting section (pp.~15--17) describes using a 2008 ACS frame and states, ``Weights larger than 7 were trimmed'' (p.~17).

\guideevidence{evidence01.png}{5.6in}{Evidence 1. 2012 guide, p. 17. Weights larger than 7 were trimmed.}

The 2010 weight seems to be capped at 7 (max 7.03, 2010 guide, p.~22), matching the yellow text above. But the 2012 codebook entry for the common content weight V103 (p.~27) reports a maximum of 20.00044. So the cap in 2012 appears not to be 7, meaning the yellow text is not correct for 2012.

\guideevidence{evidence02.png}{5.6in}{Evidence 2. 2012 guide, p. 27. The V103 codebook entry, with a maximum of 20.00044.}

Apart from this, the 2012 sampling and weighting section is word-for-word identical to the 2010 section (2010 guide, pp.~13--15), with no differences other than heading capitalization, bullet glyphs, page numbers, and one hyphenation, which suggests the section was copied over and not updated.

If that is correct, we cannot tell which parts of the section (frame, strata, panels, matching function, weighting variables) accurately describe the 2012 data and which describe the 2010 and should have been updated for 2012.

We describe the 2012 procedure as printed in the 2012 guide, because it is the only description available, and table~\ref{tab:ces_stages} gives 2012 the same entries as 2010.

\paragraph*{2. The 2014 codebook summary of the common content weight is the 2012 summary (year with possible errors: 2014)}
The 2014 guide's codebook summary of the common content weight V103 (p.~21) reports $N = 54{,}535$, mean 1, SD 1.31757, min 0.00001, max 20.00044.

This entry describes the 2012 common content weight, not the 2014 common content weight.

Every statistic in the 2014 guide matches the corresponding 2012 entry (2012 guide, p.~27); only the labels differ (e.g., ``Number of Observations'' vs.\ ``Number of Valid Observations''). This entry also cannot describe 2014 accurately. Its $N$ (54,535) is the 2012 sample size, whereas the 2014 guide states that the 2014 common content sample is 56,200 (pp.~7, 12), and its maximum (20.00044) exceeds the 2014 stated common-content trim threshold of 15 (p.~16). Table~\ref{tab:guide_v103_2012_2014} compares the two entries.

\begin{table}[!htbp]
\centering
\footnotesize
\begin{threeparttable}
\caption{The V103 codebook entries in the 2012 and 2014 guides.}
\label{tab:guide_v103_2012_2014}
\begin{tabular}{lll}
\toprule
Statistic & 2012 guide, V103 (p.~27) & 2014 guide, V103 (p.~21) \\
\midrule
$N$ & 54,535 & 54,535 \\
Mean & 1 & 1 \\
SD & 1.31757 & 1.31757 \\
Minimum & 0.00001 & 0.00001 \\
Maximum & 20.00044 & 20.00044 \\
Stated common content sample & 54,535 (p.~7) & 56,200 (pp.~7, 12) \\
Stated trimming threshold & 7 (p.~17) & 15 common content; 7 team (p.~16) \\
\bottomrule
\end{tabular}
\end{threeparttable}
\end{table}

\guideevidence{evidence03.png}{5.6in}{Evidence 3. 2014 guide, p. 21. Compare the 2012 entry in Evidence 2.}

The 2014 methods text itself appears to have been revised. Only the codebook entry seems to require updating. The error does not affect our analysis, which uses the weight in the released 2014 file (\texttt{weight}).

\paragraph*{3. The matching-pool size ``approximately 87,389'' is repeated in three guides (2014, 2016, 2018)}
The 2014 (p.~14), 2016 (p.~15), and 2018 (p.~13) guides each say the final set of completed pre-election interviews, after quality controls, was approximately 87,389, and that these were matched to the target frame.

\guideevidence{evidence04.png}{4.6in}{Evidence 4. 2014 guide, p. 14.}

\guideevidence{evidence05.png}{4.6in}{Evidence 5. 2016 guide, p. 15.}

\guideevidence{evidence06.png}{4.6in}{Evidence 6. 2018 guide, p. 13.}

An identical pool size in three different years is unlikely, which suggests the figure was carried forward.

We report the figure only for 2014, where it first appears, and treat the 2016 and 2018 figures as incorrectly carried forward without updating.

\paragraph*{4. The 2014 and 2016 lists of variables used for matching are not the same as the variables shown in the printed distance functions (2014, 2016)}
The 2008--2016 guides each provide the distance function and, beneath it, a list of variables used for matching. In each of 2008, 2010, and 2012, the variables in the list are the same as the variables in the function. In 2014 and 2016, the variable list and the function disagree in two ways. Newsint (interest in politics) is in the list but not in the function, and voter registration is in the function (weight 15) but not in the list. So at least one of the two is wrong in each year, and the guides do not say which.

\guideevidence{evidence07.png}{5.6in}{Evidence 7. 2014 guide, p. 15. The distance function; voter registration enters with weight 15; no Newsint term.}

\guideevidence{evidence08.png}{5.6in}{Evidence 8. 2014 guide, p. 15. The listed matching variables; Newsint listed, registration not; Pid5b described as ``from July 2012''.}

We think the function is more likely to be correct. Between 2014 and 2016, the function was revised: five cross-product terms were dropped (voter registration $\times$ gender, race $\times$ age, race $\times$ education, ideology $\times$ gender, gender $\times$ age), and some weights changed (the weight for age changed from 3 to 25; the weight for education indicators changed from 25 to 10 for each one; the weight for ``took both waves'' changed from 2 to 25). Yet the 2016 list of variables is word-for-word the same as the 2014 list of variables, including the description of Pid5b as baseline party ID ``from July 2012'' (2016 guide, p.~15). This suggests the function was updated for 2016 and the list was not.

Because we cannot verify which source is correct, we make the conservative choice and count both political interest and voter registration as matching variables in 2014 and 2016 (the union of the list and the function).

% --------------------------------------------------------------------------
\subsubsection*{Unclear Variables and Targets Used to Compute the Weights}

The guides describe weighting in two steps. In the first step, the matched sample is weighted to resemble the sampling frame, using propensity scores in 2008--2012 and entropy balancing in 2014--2022. Entropy balancing chooses weights so that the weighted sample reproduces specified features of the frame, such as the proportion of women or the proportion with a college degree. These features are called ``moment conditions.'' A variable affects the first-step weights only if it appears among the moment conditions.

In the second step, called ``post-stratification'' in most guides and ``raking'' in 2022, the weights are adjusted so that the sample matches population distributions on a second list of variables. This requires a population figure (a ``target'') for each of those variables.

\paragraph*{5. Hispanic origin (and, in 2022, turnout) is listed as a weighting target but not as a moment condition (2018--2022)}
The 2018, 2020, and 2022 guides each describe the first weighting step in two consecutive sentences. The first sentence lists the frame distributions the sample was weighted to match. The second lists the moment conditions. Hispanic origin appears in the first list but not the second in all three years (2018 guide, p.~14; 2020 guide, p.~16; 2022 guide, pp.~15--16). In 2022, the same is true of turnout, which is also absent from the raking list (Evidence 9).

\guideevidence{evidence09.png}{5.6in}{Evidence 9. 2022 guide, p. 16. The first sentence begins at the foot of p. 15. Hispanic origin and turnout appear among the frame distributions but not among the moment conditions. The raking list follows.}

Because entropy balancing adjusts only for the moment conditions, one of the two sentences must be wrong about these variables. We cannot tell whether the sample was actually weighted on Hispanic origin in 2018--2022, or on turnout in 2022. We assume the first sentence is correct, so that Hispanic origin (2018--2022) and turnout (2022) were used in the first weighting step; table~\ref{tab:ces_stages} classifies them this way.

\paragraph*{6. The interaction terms in the moment conditions are not specified (2014--2022)}
The 2014--2022 guides say the moment conditions included a listed set of variables ``plus their interactions'' (e.g., 2014 guide, p.~16; 2020 guide, p.~16). They do not say which interactions were used, or whether the weights matched only the proportions in each category (first moments) or also higher moments.

We assume first-moment constraints on main-effect and interaction indicators, following \textcite{hainmueller_entropy_2012}. Table~\ref{tab:ces_stages} marks ``interactions among the above'' as used in stage-one weighting for 2014--2022, without naming the specific interaction terms.

\paragraph*{7. Turnout is added to the 2014 frame but not used in any later step (2014)}
The 2014 guide says that turnout data from the November 2010 and November 2008 CPS were added to the frame (p.~14; Evidence 10). But turnout does not appear in any later step (matching on pp.~14--15, weighting on p.~16, or post-stratification on p.~16). This suggests the turnout sentence may have been carried forward from the 2008--2012 guides, when turnout was used in the propensity-score weighting. If so, it raises the possibility that the 2014 frame was described incorrectly.

\guideevidence{evidence10.png}{5.6in}{Evidence 10. 2014 guide, p. 14. Two turnout variables attached to the frame; target strata include registration, not turnout.}

We assume turnout was carried on the frame and simply unused. We record turnout as documented non-use at every stage in 2014, and in 2016--2020.

\paragraph*{8. The offices in the statewide-races adjustment are not named (2014, 2016)}
In 2014 and 2016, the guides say that the common content weights were post-stratified across states and statewide political races, without saying which races. Brian Schaffner's research assistant at Tufts confirmed to us by email that these were the governor and US Senate races, but this is not in the documentation.

We treat the 2014 and 2016 adjustment as using governor and US Senate results, as table~\ref{tab:ces_stages} shows, consistent with the 2018--2022 guides.

\paragraph*{9. No source is given for the born-again target (2018--2022)}
Born-again status is a post-stratification variable in 2018--2022 (2018 guide, p.~14; 2020 guide, p.~16; 2022 guide, p.~16; Evidence 9), but the guides do not say where its population target came from. It cannot come from the ACS, which does not ask about religion. And unlike the 2014 and 2016 guides, which say religion was attached to the frame from the 2007 Pew Religious Landscape Survey (2014 guide, p.~14; 2016 guide, p.~14), the 2018--2022 guides do not describe attaching religion from any source.

We count born-again status as a post-stratification variable in 2018--2022, with its target unknown.

\paragraph*{10. The guides are unclear about which cases were weighted to the frame (2018--2022)}
The 2018, 2020, and 2022 guides describe the first weighting step as weighting the ``completed cases'' to the frame (2018 guide, p.~14; 2020 guide, p.~16; 2022 guide, pp.~15--16). The preceding sentence in each guide describes the ``matched cases'' and the frame being combined. All completed interviews and the matched sample are different sets of respondents. The guides are therefore unclear about which set was weighted. In 2008--2016 the guides describe weighting the matched cases.

We assume that, as in earlier years, the matched cases were weighted.

% --------------------------------------------------------------------------
\subsubsection*{Steps That Are Not Documented}

The issues below are steps in the sampling and weighting process that the guides do not describe, or describe only in part, for the years listed.

\paragraph*{11. Target-sample strata (2018--2022)}
The 2008--2016 guides state that the target sample was drawn from the frame by stratified random sampling on age, race, gender, and education. Voter registration was added as a stratification variable in 2014 and also appears in 2016 (2014 guide, p.~14; 2016 guide, p.~14).

The 2018, 2020, and 2022 guides do not describe the target-sample strata, and do not state that the target sample was drawn by stratified sampling in those years.

We mark target strata as undocumented for 2018--2022 and state that the strata are unknown.

\paragraph*{12. How the two frames were combined (2016, 2018, 2020)}
The 2016 (p.~14), 2018 (p.~13), and 2020 (p.~15) guides each say that YouGov ``employed a combination of two frames.'' None of the three says how the two frames were combined. The 2016 guide describes both frames, one built from the 2012 ACS and the other from the 2010 ACS. The 2016 guide's description of its second frame, drawn from the 2010 ACS, repeats the 2014 guide's frame description, which suggests that it was carried forward from the earlier guide. The 2018 and 2020 guides describe only one of the two frames. In contrast, the 2022 guide describes a single modeled frame.

\guideevidence{evidence11.png}{5.6in}{Evidence 11. 2016 guide, p. 14. Two frames, the first from the 2012 ACS and the second from the 2010 ACS.}

\guideevidence{evidence12.png}{5.6in}{Evidence 12. 2020 guide, p. 15. The 2018 guide, p. 13, opens the same way.}

As a result, it is unclear how the two frames were combined into the frame that the sample was matched and weighted to, and, for 2018 and 2020, what the second frame was.

For 2016, we report both frames as the guide describes them (section E, table~\ref{tab:frame_sources_by_year}). For 2018 and 2020, we report the one frame each guide describes. For none of these years do we say how the frames were combined.

\paragraph*{13. How the 2022 modeled frame was built (2022)}
The 2022 guide says YouGov used a politically representative ``modeled frame'' based on the 2019 ACS, ``with modeled political variables, such as voter registration and 2020 Presidential vote,'' drawing on voter files, the 2020 CPS, the 2020 exit poll, and the 2020 CES (p.~15). The guide does not indicate how these variables were modeled or what makes the frame ``politically representative.''

We describe the 2022 frame as based on the 2019 ACS but do not attempt to describe the modeling process.

\paragraph*{14. When cases were selected in 2014 and 2016 (2014, 2016)}
The 2014 and 2016 guides say that completed pre-election interviews were matched to the target sample. But the printed distance functions include a term for whether the respondent took both waves, with weight 2 in 2014 (2014 guide, p.~15) and 25 in 2016 (2016 guide, [p.~?]). A respondent's completion of the post-election wave is not known until after the election. Either selection occurred after the post-election wave, or the function is misprinted.

\paragraph*{15. CPS registration data in the 2018 and 2020 frames (2018, 2020)}
The 2018 guide says voter registration on the frame was matched from ``the 2018 CPS'' (p.~13), and the 2020 guide says ``the 2018 and 2020 Current Population Survey (CPS)'' (p.~15; Evidence 12). The November CPS Voting and Registration Supplement for 2018 was released on April 23, 2019, and the one for 2020 on April 29, 2021. Both dates are after the election and after the target sample would have been drawn. Unless ``CPS'' refers to something other than the November supplement, these frames cite data that were not available when the frames were built. The descriptions may instead refer to data attached at a later stage; for example, the 2020 vote-validation weights, which were released later, also cite the 2020 CPS.

We report the frame sources as the guides state them (table~\ref{tab:frame_sources_by_year}) and do not use the frame's registration source in any analysis.

\paragraph*{16. The distance metric is mislabeled (2008--2016)}
The 2008--2016 guides call the matching metric a ``weighted Euclidean distance.'' The functions they print are weighted sums of absolute differences and mismatch indicators. They have no squared terms and no square root, so they are not Euclidean distances. The main text describes the metric in the form the guides print it (``Selecting Cases from the Pool of Completed Interviews'').

We follow the printed functions.

\paragraph*{17. How the post-election weights were computed is not described (2016--2022)}
The 2016 file includes pre-election and post-election weights, each in versions computed before and after vote validation. From 2018, the files include \texttt{commonweight} for all respondents and \texttt{commonpostweight} for respondents who completed both questionnaires. They also include weights for respondents matched to a validated registration record. The guides describe the target population of each weight but not how the post-election weights were computed.

We use the post-election weight for every item in 2016--2022 (Appendix~A, table~\ref{tab:disclosure_by_year}).

% \listoffigures
% \listoftables


\end{document}
